\documentclass{article}
\usepackage[margin=0.8in]{geometry}
\usepackage{amsmath,amssymb,amsthm,mathtools}
\usepackage{mathrsfs,bm,bbm}
\usepackage{graphicx,booktabs,array,multirow,tabularx,longtable}
\usepackage{enumitem}
\usepackage{xcolor}
\usepackage{algorithm,algpseudocode}
\usepackage{subcaption,tikz}
\usepackage{siunitx,cancel,accents}
\usepackage{microtype}
\usepackage[numbers,sort&compress]{natbib}
\usepackage[colorlinks=true,linkcolor=blue,citecolor=blue,urlcolor=blue]{hyperref}
\usepackage{cleveref}
\setlength{\emergencystretch}{2em}
\newtheorem{assumption}{Assumption}
\newtheorem{definition}{Definition}
\newtheorem{theorem}{Theorem}
\newtheorem{lemma}{Lemma}
\newtheorem{proposition}{Proposition}
\newtheorem{openendedquestion}{Open-ended Question}
\newtheorem{conjecture}{Conjecture}
\newcommand{\coreref}[1]{\ref{#1}}

\begin{document}

\sloppy

\section*{pid\_\allowbreak{}areal\_\allowbreak{}rd\_\allowbreak{}transport\_\allowbreak{}modulus}

\noindent\textit{Specialization:} v1\par

\paragraph{TL;DR.} Fixed geographic cells reveal integrals of latent response surfaces, not their boundary traces. Under maintained known polygonal score laws, bounded outcomes, and the same \([0,1]\)-bounded isotropic Euclidean-Lipschitz class used armwise with separately calibrated radii, the weighted boundary effect has an attained sharp interval whose support values equal clipping-aware partial-transport costs inside an unrestricted exact-moment multiplier dual. Closed-form rectangle, nonuniform-density, and weighted-triangle laws expose when fixed resolution is vacuous, informative, or sign identifying. Deterministically projected atomic outer programs and exact-moment piecewise-affine inner programs converge with quantitative relative-interior gaps. Affine-hull-projected Hoeffding regions give simultaneous finite-sample coverage and Hausdorff consistency without ambient full rank. On rational geometry and relative-interior means, the terminating exact-arithmetic \(\mathsf{VPA}\) routine returns exact continuous-moment primal witnesses and arbitrarily tight outward endpoint enclosures.

\section{Project justification}

\paragraph{Gap.} Boundary-discontinuity methods observe units approaching the assignment boundary, whereas many historical and administrative applications expose only fixed areal means. Existing tight smoothness programs do not invert these area integrals, and existing unequal-mass transport or moment hierarchies do not combine clipping, exact cell moments, causal endpoint lifting, and finite stopping certificates.

\paragraph{Niche.} The observation operator is finite-dimensional but the target is singular with respect to every cell law. This creates a clipped moment-fiber geometry in which primal extrema exist universally, outer multipliers can diverge on exposed faces, and exact continuous primal recovery becomes possible locally only through relative-interior slack.

\paragraph{Fill.} The project characterizes the exact causal set and its transport dual, proves the multiplier-attainment frontier and analytic planar formulas, derives convergent deterministic inner and outer programs, and gives affine-hull-projected whole-set inference without an ambient-rank condition. Its certified computational contribution is a terminating exact-arithmetic primal reconstruction and dual-enclosure algorithm for rational polygonal geometry, rational piecewise score laws and inputs, and relative-interior means. The maintained response class remains the \([0,1]\)-bounded isotropic Euclidean-Lipschitz class applied armwise with separately calibrated radii; novelty lies in the fixed-areal operator, weighted boundary target, inversion, and certification.

\section{Related work}

Qiu and Stoye's equation \(2.1\) imposes Euclidean Lipschitz smoothness, equation \(2.2\) gives the tight bounded isotropic boundary-to-policy program, and equation \(2.3\) deliberately weakens it to one coordinate. Their open computational program informs the smoothness geometry but is not executed here: the direction of extrapolation and the fixed-areal observation operator are different. Piccoli and Rossi's Definition 12 and Theorem 13 identify \(W_1^{1,1}\) with the bounded-Lipschitz flat dual, while their Definition 3 and Proposition 4 give the attained equal-mass submeasure representation. Those two cited leaves supply only the classical transport substrate: the clipping rescaling, exact areal-moment fiber, unrestricted multiplier dual, relative-interior multiplier attainment, and boundary nonattainment witness are derived here. Cattaneo, Titiunik, and Yu define weighted boundary estimands and develop unit-level pre- and post-aggregation inference, which fixes the causal target but not its areal mean-information set. Mula and Nouy provide the moment-hierarchy template for value-side approximation, while Baldi and Mourrain delimit when finite exact moment extraction follows from flat truncation. Moon shows that aggregate-data sharp sets and optimization are established themes for finite-support covariates, making the continuous singular-boundary and global-Lipschitz structure essential to novelty. Chernozhukov, Hong, and Tamer provide the general set-inversion benchmark; the present inferential claim uses exact finite-sample bounded-mean coverage and then the sharp endpoint map. Dell supplies the concrete bounded geographic-RD outcome for which a fixed-resolution sign certificate has direct substantive meaning.

\section{Setup}

Target (estimand / structural parameter): \(\tau_B=\int_B\{f_1(x)-f_0(x)\}\,d\nu(x)\).
Primary functional / estimator / diagnostic: \(I(m_0,m_1)=[L_1(m_1)-U_0(m_0),\ U_1(m_1)-L_0(m_0)]\), with \(U_d(m_d)=\inf_{\lambda_d\in\mathbb R^{|J_d|}}\{\lambda_d^\top m_d+H_{M_d}(\nu-\sum_j\lambda_{dj}\mu_{dj})\}\) and \(L_d(m_d)=1-U_d(\mathbf 1-m_d)\).
\begin{itemize}
  \item \texttt{d} --- \texttt{arm index}; \(\{0,1\}\); \texttt{Indexes control and treatment potential-\allowbreak{}outcome surfaces.}
  \item \texttt{S} --- \texttt{score domain}; \(\mathbb R^2\); \texttt{Common spatial domain.}
  \par\smallskip\noindent\emph{Definition.} \(S\) is the compact polygonal geographic score support.
  \item \texttt{X} --- \texttt{geographic score}; \(S\); \texttt{Assignment variable.}
  \par\smallskip\noindent\emph{Definition.} \(X\) records the two-dimensional location score.
  \item \texttt{Y\_\allowbreak{}d} --- \texttt{potential outcome}; \([0,1]\); \texttt{Primitive causal outcome.}
  \par\smallskip\noindent\emph{Definition.} \(Y(d)\) is the potential outcome under arm \(d\).
  \item \texttt{B} --- \texttt{assignment boundary}; closed polygonal subset of \(S\); \texttt{Support of the causal target.}
  \par\smallskip\noindent\emph{Definition.} \(B\) separates the two assignment sides of \(S\).
  \item \texttt{S\_\allowbreak{}d} --- \texttt{assignment side}; closed subset of \(S\); \texttt{Support of arm-\allowbreak{}specific observed cells.}
  \par\smallskip\noindent\emph{Definition.} \(S_d\) is the closure of the side assigned arm \(d\), with \(S_0\cap S_1=B\).
  \item \texttt{J\_\allowbreak{}d} --- \texttt{finite cell index set}; finite subsets of \(\mathbb N\); \texttt{Dimension of the arm-\allowbreak{}specific information vector.}
  \par\smallskip\noindent\emph{Definition.} \(J_d\) indexes the observed areal cells on side \(d\).
  \item \texttt{C\_\allowbreak{}dj} --- \texttt{areal cell}; \texttt{compact rational polygons}; \texttt{Spatial aggregation unit.}
  \par\smallskip\noindent\emph{Definition.} \(C_{dj}\subset S_d\) for \(j\in J_d\).
  \item \texttt{mu\_\allowbreak{}dj} --- \texttt{within-\allowbreak{}cell score law}; probability measures on \(S\); \texttt{Areal observation functional.}
  \par\smallskip\noindent\emph{Definition.} \(\mu_{dj}\) is the conditional law of \(X\) in observed cell \(C_{dj}\).
  \item \texttt{w} --- \texttt{boundary weight}; nonnegative piecewise-affine functions on \(B\); \texttt{Weights heterogeneous effects along the boundary.}
  \par\smallskip\noindent\emph{Definition.} \(w\) is the known density of the target measure relative to boundary length.
  \item \texttt{nu} --- \texttt{known boundary law}; probability measures on \(B\); \texttt{Weights the boundary-\allowbreak{}average causal effect.}
  \par\smallskip\noindent\emph{Definition.} \(d\nu=w\,d\mathcal H^1\) on \(B\).
  \item \texttt{M\_\allowbreak{}d} --- \texttt{Lipschitz radius}; \([0,\infty)\); \texttt{Arm-\allowbreak{}specific smoothness sensitivity input.}
  \par\smallskip\noindent\emph{Definition.} \(M_d\) is the separately selected maintained Lipschitz radius for arm \(d\).
  \item \texttt{g\_\allowbreak{}d} --- \texttt{generic response surface}; \(C(S)\); \texttt{Bound variable in the response-\allowbreak{}surface class.}
  \par\smallskip\noindent\emph{Definition.} \(g_d:S\to\mathbb R\) is a generic class member for arm \(d\).
  \item \texttt{F\_\allowbreak{}d} --- \texttt{response-\allowbreak{}surface class}; subsets of \(C(S)\); \texttt{Primitive structural class.}
  \par\smallskip\noindent\emph{Definition.} \(\mathcal F_d\) is the clipped \(M_d\)-Lipschitz class defined in \(\mathrm{def:surface-class}\).
  \item \texttt{f\_\allowbreak{}d} --- \texttt{conditional mean potential-\allowbreak{}outcome surface}; \(\mathcal F_d\); \texttt{Latent arm-\allowbreak{}specific causal response.}
  \par\smallskip\noindent\emph{Definition.} \(f_d(x)=\mathbb E[Y(d)\mid X=x]\).
  \item \texttt{A\_\allowbreak{}d} --- \texttt{areal moment operator}; maps \(\mathcal F_d\) to \(\mathbb R^{|J_d|}\); \texttt{Maps a surface to observed cell means.}
  \item \texttt{m\_\allowbreak{}d} --- \texttt{observed mean vector}; \(\mathbb R^{|J_d|}\); Population information for arm \(d\).
  \par\smallskip\noindent\emph{Definition.} \(m_d=A_df_d\).
  \item \texttt{K\_\allowbreak{}d} --- \texttt{moment body}; compact convex subsets of \(\mathbb R^{|J_d|}\); \texttt{Set of compatible mean vectors.}
  \par\smallskip\noindent\emph{Definition.} \(K_d=A_d(\mathcal F_d)\).
  \item \texttt{tau\_\allowbreak{}B} --- \texttt{causal estimand}; \([-1,1]\); \texttt{Weighted boundary-\allowbreak{}average treatment effect.}
  \par\smallskip\noindent\emph{Definition.} \(\tau_B=\int_B(f_1-f_0)\,d\nu\).
  \item \texttt{U\_\allowbreak{}d} --- \texttt{upper endpoint map}; maps \(K_d\) to \([0,1]\); Sharp upper boundary mean for arm \(d\).
  \par\smallskip\noindent\emph{Definition.} \(U_d(m_d)=\max\{\int_B g_d\,d\nu:g_d\in\mathcal F_d,\ A_dg_d=m_d\}\).
  \item \texttt{L\_\allowbreak{}d} --- \texttt{lower endpoint map}; maps \(K_d\) to \([0,1]\); Sharp lower boundary mean for arm \(d\).
  \par\smallskip\noindent\emph{Definition.} \(L_d(m_d)=\min\{\int_B g_d\,d\nu:g_d\in\mathcal F_d,\ A_dg_d=m_d\}\).
  \item \texttt{I\_\allowbreak{}m} --- \texttt{identified set}; compact intervals in \(\mathbb R\); Sharp mean-information set for \(\tau_B\).
  \par\smallskip\noindent\emph{Definition.} \(I(m_0,m_1)=[L_1(m_1)-U_0(m_0),\ U_1(m_1)-L_0(m_0)]\).
  \item \texttt{sigma} --- \texttt{finite signed measure}; \(\mathcal M(S)\); \texttt{Residual measure in the support problem.}
  \par\smallskip\noindent\emph{Definition.} \(\sigma=\sigma^+-\sigma^-\) denotes its Jordan decomposition.
  \item \texttt{sigma\_\allowbreak{}plus} --- \texttt{positive Jordan part}; nonnegative finite measures on \(S\); \texttt{Positive supply in partial transport.}
  \par\smallskip\noindent\emph{Definition.} \(\sigma^+\) is the positive part of \(\sigma\).
  \item \texttt{sigma\_\allowbreak{}minus} --- \texttt{negative Jordan part}; nonnegative finite measures on \(S\); \texttt{Negative demand in partial transport.}
  \par\smallskip\noindent\emph{Definition.} \(\sigma^-\) is the negative part of \(\sigma\).
  \item \texttt{H\_\allowbreak{}Md} --- \texttt{clipped Lipschitz support value}; maps signed measures to \(\mathbb R\); \texttt{Inner support cost in the endpoint dual.}
  \par\smallskip\noindent\emph{Definition.} \(H_{M_d}(\sigma)=\max_{g_d\in\mathcal F_d}\int_S g_d\,d\sigma\).
  \item \texttt{Pi\_\allowbreak{}le\_\allowbreak{}sigma} --- \texttt{partial-\allowbreak{}coupling set}; nonnegative measures on \(S\times S\); \texttt{Feasible subtransport plans.}
  \par\smallskip\noindent\emph{Definition.} \(\Pi_{\le}(\sigma)=\{\pi\ge0:\pi_1\le\sigma^+,\ \pi_2\le\sigma^-\}\).
  \item \texttt{lambda\_\allowbreak{}d} --- \texttt{moment multiplier}; \(\mathbb R^{|J_d|}\); \texttt{Dualizes exact areal moments.}
  \par\smallskip\noindent\emph{Definition.} \(\lambda_d=(\lambda_{dj})_{j\in J_d}\) is unrestricted.
  \item \texttt{D\_\allowbreak{}d} --- \texttt{dual objective}; maps \(\mathbb R^{|J_d|}\times K_d\) to \(\mathbb R\); \texttt{Finite-\allowbreak{}multiplier endpoint bound.}
  \par\smallskip\noindent\emph{Definition.} \(D_d(\lambda_d;m_d)=\lambda_d^\top m_d+H_{M_d}(\nu-\sum_{j\in J_d}\lambda_{dj}\mu_{dj})\).
  \item \texttt{E\_\allowbreak{}d} --- \texttt{relative tangent space}; linear subspaces of \(\mathbb R^{|J_d|}\); \texttt{Affine geometry of the moment body.}
  \par\smallskip\noindent\emph{Definition.} \(E_d=\operatorname{span}(K_d-K_d)\).
  \item \texttt{r\_\allowbreak{}d} --- \texttt{relative-\allowbreak{}interior radius}; \([0,\infty)\); \texttt{Quantifies local feasible moment slack.}
  \par\smallskip\noindent\emph{Definition.} \(r_d=r_d(m_d)=\sup\{r\ge0:m_d+r\mathbb B_{E_d}\subseteq K_d\}\).
  \item \texttt{h} --- \texttt{mesh diameter}; \((0,\infty)\); \texttt{Resolution of a conforming refinement.}
  \item \texttt{T\_\allowbreak{}h} --- \texttt{conforming triangulation}; finite rational triangulations of \(S\); \texttt{Finite geometric discretization.}
  \par\smallskip\noindent\emph{Definition.} Fix one rational triangulation \(\mathcal T_0\) conforming to \(B\), all cell boundaries, and all density and boundary-weight knots. Recursively form \(\mathcal T_{k+1}\) by joining the three rational edge midpoints of every triangle of \(\mathcal T_k\), and write \(\mathcal T_h=\mathcal T_k\) when \(h=\max_{T\in\mathcal T_k}\operatorname{diam}(T)\).
  \item \texttt{N\_\allowbreak{}h} --- \texttt{rational nodal set}; finite subsets of \(S\); \texttt{Common finite support for the atomic program.}
  \par\smallskip\noindent\emph{Definition.} \(\mathcal N_h\) is the vertex set of \(\mathcal T_h\), ordered lexicographically by its rational coordinates.
  \item \texttt{p\_\allowbreak{}djh} --- \texttt{deterministic cell nodal projection}; maps \(C_{dj}\) to \(\mathcal N_h\cap C_{dj}\); \texttt{Pushes each cell score law to rational mesh nodes without ambiguity.}
  \par\smallskip\noindent\emph{Definition.} \(p_{dj,h}(x)\) is the lexicographically first element of \(\operatorname*{argmin}_{z\in\mathcal N_h\cap C_{dj}}\|x-z\|_2\).
  \item \texttt{p\_\allowbreak{}Bh} --- \texttt{deterministic boundary nodal projection}; maps \(B\) to \(\mathcal N_h\cap B\); \texttt{Pushes the boundary law to rational boundary nodes without ambiguity.}
  \par\smallskip\noindent\emph{Definition.} \(p_{B,h}(x)\) is the lexicographically first element of \(\operatorname*{argmin}_{z\in\mathcal N_h\cap B}\|x-z\|_2\).
  \item \texttt{mu\_\allowbreak{}djh} --- \texttt{atomic cell law}; probability measures on \(\mathcal N_h\cap C_{dj}\); \texttt{Exact pushed-\allowbreak{}forward cell weights in the atomic moment constraints.}
  \par\smallskip\noindent\emph{Definition.} \(\mu_{dj,h}=(p_{dj,h})_\#\mu_{dj}\).
  \item \texttt{nu\_\allowbreak{}h} --- \texttt{atomic boundary law}; probability measures on \(\mathcal N_h\cap B\); \texttt{Exact pushed-\allowbreak{}forward weights in the atomic boundary objective.}
  \par\smallskip\noindent\emph{Definition.} \(\nu_h=(p_{B,h})_\#\nu\).
  \item \texttt{v\_\allowbreak{}dh} --- \texttt{nodal response vector}; \([0,1]^{|\mathcal N_h|}\); \texttt{Decision variable of the finite atomic relaxation.}
  \par\smallskip\noindent\emph{Definition.} \(v_{d,h}=(v_{d,h,z})_{z\in\mathcal N_h}\).
  \item \texttt{E\_\allowbreak{}dh} --- \texttt{deterministic clipped McShane extension}; maps feasible nodal vectors to \(\mathcal F_d\); \texttt{Converts every pairwise-\allowbreak{}Lipschitz nodal solution into a canonical continuous feasible surface.}
  \par\smallskip\noindent\emph{Definition.} \((E_{d,h}v_{d,h})(x)=\min\{1,\max\{0,\min_{z\in\mathcal N_h}[v_{d,h,z}+M_d\|x-z\|_2]\}\}\).
  \item \texttt{F\_\allowbreak{}dh} --- \texttt{piecewise-\allowbreak{}affine inner class}; finite-dimensional subsets of \(\mathcal F_d\); \texttt{Exactly feasible primal approximation class.}
  \par\smallskip\noindent\emph{Definition.} \(\mathcal F_{d,h}\) contains the continuous functions affine on each triangle of \(\mathcal T_h\) that satisfy the exact range and Euclidean Lipschitz inequalities.
  \item \texttt{epsilon\_\allowbreak{}dh} --- \texttt{uniform approximation error}; \([0,\infty)\); \texttt{Controls exact-\allowbreak{}moment inner correction.}
  \par\smallskip\noindent\emph{Definition.} \(\varepsilon_{d,h}=\sup_{g_d\in\mathcal F_d}\inf_{q\in\mathcal F_{d,h}}\|g_d-q\|_\infty\).
  \item \texttt{P\_\allowbreak{}dh} --- \texttt{inner endpoint value}; \([0,1]\); Certified exactly moment-feasible lower bound on \(U_d\).
  \par\smallskip\noindent\emph{Definition.} \(P_{d,h}(m_d)=\max\{\int_Bq\,d\nu:q\in\mathcal F_{d,h},\ A_dq=m_d\}\).
  \item \texttt{V\_\allowbreak{}dh} --- \texttt{atomic nodal endpoint value}; \([0,1]\); \texttt{Finite atomic relaxation.}
  \par\smallskip\noindent\emph{Definition.} \(V_{d,h}(m_d)\) is the maximum of \(\sum_{z\in\mathcal N_h}v_{d,h,z}\nu_h(\{z\})\) subject to \(0\le v_{d,h,z}\le1\), \(|v_{d,h,z}-v_{d,h,z'}|\le M_d\|z-z'\|_2\) for all \(z,z'\in\mathcal N_h\), and \(\left|\sum_zv_{d,h,z}\mu_{dj,h}(\{z\})-m_{dj}\right|\le M_dh\) for every \(j\in J_d\).
  \item \texttt{B\_\allowbreak{}dh} --- \texttt{outer endpoint value}; \([0,1]\); Universal atomic outer bound on \(U_d\).
  \par\smallskip\noindent\emph{Definition.} \(B_{d,h}(m_d)=\min\{1,V_{d,h}(m_d)+M_dh\}\).
  \item \texttt{N} --- \texttt{number of observed cells}; \(\mathbb N\); \texttt{Multiplicity in simultaneous concentration.}
  \par\smallskip\noindent\emph{Definition.} \(N=|J_0|+|J_1|\).
  \item \texttt{n\_\allowbreak{}dj} --- \texttt{cell sample size}; \texttt{positive integers}; \texttt{Cell-\allowbreak{}specific information size.}
  \par\smallskip\noindent\emph{Definition.} \(n_{dj}\) is fixed conditional on the sampling design.
  \item \texttt{Y\_\allowbreak{}dji} --- \texttt{observed cell outcome}; \([0,1]\); \texttt{Sampling observation.}
  \par\smallskip\noindent\emph{Definition.} \(Y_{dji}\) has mean \(m_{dj}\) for \(i=1,\ldots,n_{dj}\).
  \item \texttt{mhat\_\allowbreak{}dj} --- \texttt{sample cell mean}; \([0,1]\); \texttt{Observed sufficient statistic for confidence inversion.}
  \par\smallskip\noindent\emph{Definition.} \(\widehat m_{dj}=n_{dj}^{-1}\sum_{i=1}^{n_{dj}}Y_{dji}\).
  \item \texttt{gamma} --- \texttt{error probability}; \((0,1)\); \texttt{Controls simultaneous coverage.}
  \item \texttt{e\_\allowbreak{}dj} --- \texttt{cell concentration radius}; \((0,\infty)\); Hoeffding half-width for cell \((d,j)\).
  \par\smallskip\noindent\emph{Definition.} \(e_{dj}=\sqrt{\log(2N/\gamma)/(2n_{dj})}\).
  \item \texttt{a} --- \texttt{rectangle width}; \([0,\infty)\); \texttt{Width normal to a rectangular boundary.}
  \item \texttt{ell\_\allowbreak{}d} --- \texttt{dimensionless rectangle radius}; \([0,\infty)\); \texttt{Indexes clipping regimes in the rectangle formulas.}
  \par\smallskip\noindent\emph{Definition.} \(\ell_d=M_da\).
  \item \texttt{t} --- \texttt{candidate boundary height}; \([0,1]\); \texttt{Scalar argument of the rectangle envelope.}
  \item \texttt{A\_\allowbreak{}ell} --- \texttt{minimal rectangle mean envelope}; maps \([0,1]\) to \([0,1]\); \texttt{Inverts the complete uniform-\allowbreak{}rectangle endpoint,\allowbreak{} including the zero-\allowbreak{}radius case.}
  \par\smallskip\noindent\emph{Definition.} \(A_0(t)=t\), and for \(\ell_d>0\), \(A_{\ell_d}(t)=[t^2-(t-\ell_d)_+^2]/(2\ell_d)\).
  \item \texttt{rho\_\allowbreak{}star} --- \texttt{nonuniform rectangle density}; densities on \([0,1]\); \texttt{Exhibits density-\allowbreak{}sensitive clipping transitions.}
  \par\smallskip\noindent\emph{Definition.} \(\rho_\star(x)=3/2\) on \([0,1/2]\) and \(\rho_\star(x)=1/2\) on \((1/2,1]\).
  \item \texttt{eps} --- \texttt{requested certificate tolerance}; \((0,\infty)\); \texttt{Target outward endpoint gap.}
  \par\smallskip\noindent\emph{Definition.} \(\epsilon\) is supplied as an exactly represented rational number.
  \item \texttt{VPA} --- \texttt{verified primal-\allowbreak{}-\allowbreak{}dual reconstruction algorithm}; \texttt{exact-\allowbreak{}arithmetic procedures on the theorem's rational relative-\allowbreak{}interior input class}; \texttt{Produces exact primal witnesses and outward endpoint enclosures on rational geometry at relative-\allowbreak{}interior means.}
  \par\smallskip\noindent\emph{Definition.} \(\mathsf{VPA}\) dovetails conforming rational refinement, exact-moment rational-nodal primal witnesses, unrestricted rational multipliers, rational atomic partial-transport plans, and outward rational interval arithmetic until the requested endpoint gap is certified.
  \item \texttt{mtilde\_\allowbreak{}d} --- \texttt{affine-\allowbreak{}hull-\allowbreak{}projected sample mean}; \(\operatorname{aff}K_d\); \texttt{Centers the arm-\allowbreak{}specific confidence region within the estimable affine geometry.}
  \par\smallskip\noindent\emph{Definition.} \(\widetilde m_d=P_{\operatorname{aff}K_d}(\widehat m_d)\), where the projection is Euclidean orthogonal projection.
  \item \texttt{G\_\allowbreak{}sens} --- \texttt{application sensitivity grid}; \texttt{finite sets of score-\allowbreak{}law and smoothness specifications}; \texttt{Prevents a sign conclusion from resting on one unreported score-\allowbreak{}law or smoothness choice.}
  \par\smallskip\noindent\emph{Definition.} \(\mathcal G_{\mathrm{sens}}\) is a prespecified finite collection of pairs \(\bigl((\mu_{dj})_{d,j},(M_0,M_1)\bigr)\) constructed by the protocol in \(\mathrm{def:application-sensitivity-protocol}\).
  \item \texttt{rho\_\allowbreak{}dn} --- \texttt{arm-\allowbreak{}specific simultaneous Euclidean radius}; \([0,\infty)\); \texttt{Converts simultaneous coordinatewise Hoeffding control into an armwise Euclidean radius.}
  \par\smallskip\noindent\emph{Definition.} \(\rho_{d,n}=(\sum_{j\in J_d}e_{dj}^2)^{1/2}\).
  \item \texttt{Qtilde\_\allowbreak{}dn} --- \texttt{arm-\allowbreak{}specific projected confidence region}; compact convex subsets of \(\operatorname{aff}K_d\); \texttt{Carries finite-\allowbreak{}sample mean uncertainty in the relative affine geometry of the moment body.}
  \par\smallskip\noindent\emph{Definition.} \(\widetilde Q_{d,n}=\{q\in\operatorname{aff}K_d:\|q-\widetilde m_d\|_2\leq\rho_{d,n}\}\).
  \item \texttt{Q\_\allowbreak{}n} --- \texttt{joint affine-\allowbreak{}hull-\allowbreak{}projected moment region}; subsets of \(\operatorname{aff}K_0\times\operatorname{aff}K_1\); \texttt{Finite-\allowbreak{}sample simultaneous region for both arm-\allowbreak{}specific areal mean vectors.}
  \par\smallskip\noindent\emph{Definition.} \(Q_n=\widetilde Q_{0,n}\times\widetilde Q_{1,n}\).
  \item \texttt{C\_\allowbreak{}n} --- \texttt{whole-\allowbreak{}set confidence interval}; compact intervals in \(\mathbb R\) or the empty set; \texttt{Confidence set that contains the full identified interval.}
  \par\smallskip\noindent\emph{Definition.} \(C_n=\bigcup_{(q_0,q_1)\in Q_n\cap(K_0\times K_1)}I(q_0,q_1)\).
\end{itemize}

\section{Assumptions}

\begin{assumption}[ass:polygonal-score-support]\label{ass:polygonal-score-support}
\(S\) is a compact convex rational polygon.
\par\smallskip\noindent\emph{Standard} (Compact semialgebraic support specialized to a rational polygon, \cite{MulaNouy2024MomentSOSOT}).
\end{assumption}

\begin{assumption}[ass:polygonal-boundary]\label{ass:polygonal-boundary}
\(B\) is a finite union of rational line segments.
\par\smallskip\noindent\emph{Standard} (Piecewise-linear boundary-discontinuity geometry, \cite{CattaneoTitiunikYu2025LocationBoundary}).
\end{assumption}

\begin{assumption}[ass:boundary-separation]\label{ass:boundary-separation}
\(S_0\cap S_1=B\) and \(S_0\cup S_1=S\).
\par\smallskip\noindent\emph{Standard} (Sharp boundary assignment, \cite{CattaneoTitiunikYu2025LocationBoundary}).
\end{assumption}

\begin{assumption}[ass:cell-disjointness]\label{ass:cell-disjointness}
For each \(d\), the interiors of \(C_{dj}\), \(j\in J_d\), are pairwise disjoint.
\par\smallskip\noindent\emph{Standard} (Spatial-cell version of fixed aggregate categories, \cite{Moon2024AggregatePartialID}).
\end{assumption}

\begin{assumption}[ass:cell-coverage]\label{ass:cell-coverage}
For each \(d\), the union of \(C_{dj}\), \(j\in J_d\), covers the observed portion of \(S_d\) up to area zero.
\par\smallskip\noindent\emph{Standard} (Spatial-cell version of aggregate-category coverage, \cite{Moon2024AggregatePartialID}).
\end{assumption}

\begin{assumption}[ass:cell-law-support]\label{ass:cell-law-support}
\(\mu_{dj}(C_{dj})=1\) for every \(d\) and \(j\in J_d\).
\par\smallskip\noindent\emph{Standard} (Within-category conditional-score support, \cite{Moon2024AggregatePartialID}).
\end{assumption}

\begin{assumption}[ass:known-cell-score-laws]\label{ass:known-cell-score-laws}
Every within-cell score law \(\mu_{dj}\) is known to the analyst rather than estimated from the outcome sample.
\par\smallskip\noindent\emph{Novel.} This maintained input is credible when administrative population registers or complete geocoded sampling frames determine within-cell location weights independently of the outcome sample; applications without such records must vary the score laws in the prespecified sensitivity grid.
\end{assumption}

\begin{assumption}[ass:boundary-null-cells]\label{ass:boundary-null-cells}
\(\mu_{dj}(B)=0\) for every \(d\) and \(j\in J_d\).
\par\smallskip\noindent\emph{Standard} (Continuous boundary score law, \cite{CattaneoTitiunikYu2025LocationBoundary}).
\end{assumption}

\begin{assumption}[ass:cell-density]\label{ass:cell-density}
Each \(\mu_{dj}\) has a rational piecewise-constant Lebesgue density on a conforming rational triangulation.
\par\smallskip\noindent\emph{Novel.} Rational piecewise-constant densities cover uniform-within-polygon records and histogram weights exported from administrative population rasters, while permitting exact cell integration on a rational refinement.
\end{assumption}

\begin{assumption}[ass:positive-cell-density]\label{ass:positive-cell-density}
The density of each \(\mu_{dj}\) is bounded above and bounded away from zero on \(C_{dj}\).
\par\smallskip\noindent\emph{Standard} (Positive bounded score density, \cite{CattaneoTitiunikYu2025LocationBoundary}).
\end{assumption}

\begin{assumption}[ass:boundary-weight-shape]\label{ass:boundary-weight-shape}
\(w\) belongs to the cone of nonnegative piecewise-affine functions on \(B\).
\par\smallskip\noindent\emph{Standard} (Known weighted boundary aggregation, \cite{CattaneoTitiunikYu2025AggregationBoundary}).
\end{assumption}

\begin{assumption}[ass:boundary-weight-normalization]\label{ass:boundary-weight-normalization}
\(\nu(B)=1\).
\par\smallskip\noindent\emph{Standard} (Normalized boundary aggregation, \cite{CattaneoTitiunikYu2025AggregationBoundary}).
\end{assumption}

\begin{assumption}[ass:surface-range]\label{ass:surface-range}
\(0\le g_d(x)\le1\) for every \(x\in S\) and \(d\in\{0,1\}\).
\par\smallskip\noindent\emph{Standard} (Bounded conditional mean response, \cite{QiuStoye2026BoundaryDecision}).
\end{assumption}

\begin{assumption}[ass:surface-lipschitz]\label{ass:surface-lipschitz}
\(|g_d(x)-g_d(y)|\le M_d\|x-y\|_2\) for every \(x,y\in S\) and \(d\in\{0,1\}\).
\par\smallskip\noindent\emph{Novel.} Separate maintained radii allow treatment and control response surfaces to have different spatial regularity; applications select each radius from outcome-scale scientific bounds and report the full prespecified sensitivity frontier rather than treating either radius as known from boundary-RD theory.
\end{assumption}

\begin{assumption}[ass:mean-compatibility]\label{ass:mean-compatibility}
\(m_d\in K_d\) for each \(d\in\{0,1\}\).
\par\smallskip\noindent\emph{Standard} (Compatibility with a sharp identified set, \cite{BeresteanuMolchanovMolinari2011RandomSets}).
\end{assumption}

\begin{assumption}[ass:relative-interior]\label{ass:relative-interior}
\(r_d(m_d)>0\) for each arm to which the relative-interior computational claim is applied.
\par\smallskip\noindent\emph{Standard} (Strict relative-interior constraint qualification, \cite{KaidoMolinariStoye2022ConstraintQualifications}).
\end{assumption}

\begin{assumption}[ass:cell-sample-independence]\label{ass:cell-sample-independence}
Conditional on the fixed counts \(n_{dj}\), the variables \(Y_{dji}\) are mutually independent.
\par\smallskip\noindent\emph{Standard} (Independent bounded summands, \cite{Hoeffding1963ProbabilityInequalities}).
\end{assumption}

\section{Definitions}

\begin{definition}[def:surface-class]\label{def:surface-class}
\par\smallskip\noindent\emph{Defined object.} \texttt{Clipped Lipschitz response class}
\par\smallskip\noindent\emph{Construction.} \(\mathcal F_d=\{g_d\in C(S):0\le g_d(x)\le1\ \text{for all }x\in S,\ |g_d(x)-g_d(y)|\le M_d\|x-y\|_2\ \text{for all }x,y\in S\}\), for every \(d\in\{0,1\}\).
\par\smallskip\noindent\emph{Member properties.} \texttt{ass:\allowbreak{}surface-\allowbreak{}range}; \texttt{ass:\allowbreak{}surface-\allowbreak{}lipschitz}.
\end{definition}

\begin{definition}[def:moment-map]\label{def:moment-map}
\par\smallskip\noindent\emph{Defined object.} \texttt{Areal moment map}
\par\smallskip\noindent\emph{Construction.} \(A_d:\mathcal F_d\to\mathbb R^{|J_d|}\) is given by \((A_dg_d)_j=\int_Sg_d\,d\mu_{dj}\) for every \(j\in J_d\) and \(d\in\{0,1\}\).
\par\smallskip\noindent\emph{Inputs.} \texttt{g\_\allowbreak{}d}.
\end{definition}

\begin{definition}[def:moment-body]\label{def:moment-body}
\par\smallskip\noindent\emph{Defined object.} \texttt{Compatible moment body}
\par\smallskip\noindent\emph{Construction.} \(K_d=A_d(\mathcal F_d)\) for every \(d\in\{0,1\}\).
\par\smallskip\noindent\emph{Inputs.} \texttt{A\_\allowbreak{}d}; \texttt{F\_\allowbreak{}d}.
\end{definition}

\begin{definition}[def:arm-endpoints]\label{def:arm-endpoints}
\par\smallskip\noindent\emph{Defined object.} \texttt{Sharp arm-\allowbreak{}specific boundary endpoints}
\par\smallskip\noindent\emph{Construction.} \(U_d(m_d)=\max_{g_d\in\mathcal F_d:A_dg_d=m_d}\int_Bg_d\,d\nu\) and \(L_d(m_d)=\min_{g_d\in\mathcal F_d:A_dg_d=m_d}\int_Bg_d\,d\nu\).
\par\smallskip\noindent\emph{Inputs.} \texttt{m\_\allowbreak{}d}.
\end{definition}

\begin{definition}[def:effect-set]\label{def:effect-set}
\par\smallskip\noindent\emph{Defined object.} \texttt{Mean-\allowbreak{}information effect set}
\par\smallskip\noindent\emph{Construction.} \(I(m_0,m_1)=\{\int_B(g_1-g_0)\,d\nu:g_d\in\mathcal F_d,\ A_dg_d=m_d,\ d\in\{0,1\}\}\).
\par\smallskip\noindent\emph{Inputs.} \texttt{m\_\allowbreak{}d}.
\end{definition}

\begin{definition}[def:support-cost]\label{def:support-cost}
\par\smallskip\noindent\emph{Defined object.} \texttt{Clipped Lipschitz support cost}
\par\smallskip\noindent\emph{Construction.} \(H_{M_d}(\sigma)=\max_{g_d\in\mathcal F_d}\int_Sg_d\,d\sigma\).
\par\smallskip\noindent\emph{Inputs.} \texttt{sigma}; \texttt{M\_\allowbreak{}d}.
\end{definition}

\begin{definition}[def:partial-couplings]\label{def:partial-couplings}
\par\smallskip\noindent\emph{Defined object.} \texttt{Partial transport plans}
\par\smallskip\noindent\emph{Construction.} \(\Pi_{\le}(\sigma)=\{\pi\in\mathcal M_+(S\times S):\pi_1\le\sigma^+,\ \pi_2\le\sigma^-\}\).
\par\smallskip\noindent\emph{Inputs.} \texttt{sigma}.
\end{definition}

\begin{definition}[def:dual-objective]\label{def:dual-objective}
\par\smallskip\noindent\emph{Defined object.} \texttt{Unrestricted moment-\allowbreak{}multiplier objective}
\par\smallskip\noindent\emph{Construction.} \(D_d(\lambda_d;m_d)=\lambda_d^\top m_d+H_{M_d}(\nu-\sum_{j\in J_d}\lambda_{dj}\mu_{dj})\), with \(\lambda_d\in\mathbb R^{|J_d|}\).
\par\smallskip\noindent\emph{Inputs.} \texttt{lambda\_\allowbreak{}d}; \texttt{m\_\allowbreak{}d}.
\end{definition}

\begin{definition}[def:relative-interior-radius]\label{def:relative-interior-radius}
\par\smallskip\noindent\emph{Defined object.} \texttt{Relative moment slack}
\par\smallskip\noindent\emph{Construction.} \(E_d=\operatorname{span}(K_d-K_d)\) and \(r_d(m_d)=\sup\{r\ge0:m_d+r\mathbb B_{E_d}\subseteq K_d\}\).
\par\smallskip\noindent\emph{Inputs.} \texttt{m\_\allowbreak{}d}; \texttt{K\_\allowbreak{}d}.
\end{definition}

\begin{definition}[def:piecewise-affine-inner]\label{def:piecewise-affine-inner}
\par\smallskip\noindent\emph{Defined object.} \texttt{Exact-\allowbreak{}moment piecewise-\allowbreak{}affine program}
\par\smallskip\noindent\emph{Construction.} On the deterministic nested midpoint refinement \(\mathcal T_h\), let \(\mathcal F_{d,h}=\{q\in C(S):q|_T\text{ is affine},\ 0\le q(z)\le1\text{ for every }z\in\mathcal N_h,\ \|\nabla(q|_T)\|_2\le M_d\text{ for every }T\in\mathcal T_h\}\) and set \(P_{d,h}(m_d)=\max\{\int_Bq\,d\nu:q\in\mathcal F_{d,h},\ A_dq=m_d\}\).
\par\smallskip\noindent\emph{Inputs.} \texttt{T\_\allowbreak{}h}; \texttt{m\_\allowbreak{}d}.
\end{definition}

\begin{definition}[def:atomic-outer]\label{def:atomic-outer}
\par\smallskip\noindent\emph{Defined object.} \texttt{Atomic outer program}
\par\smallskip\noindent\emph{Construction.} For each \(C_{dj}\) and for \(B\), project a point to its nearest node in \(\mathcal N_h\) lying in the same set, breaking distance ties by the lexicographic order of rational node coordinates. Thus \(p_{dj,h}\) and \(p_{B,h}\) are single-valued and satisfy \(\|x-p_{dj,h}(x)\|_2\le h\) and \(\|x-p_{B,h}(x)\|_2\le h\). Set \(\mu_{dj,h}=(p_{dj,h})_\#\mu_{dj}\) and \(\nu_h=(p_{B,h})_\#\nu\). Define \(V_{d,h}(m_d)\) as the maximum of \(\sum_{z\in\mathcal N_h}v_{d,h,z}\nu_h(\{z\})\) over \(v_{d,h}\) satisfying \(0\le v_{d,h,z}\le1\), \(|v_{d,h,z}-v_{d,h,z'}|\le M_d\|z-z'\|_2\) for every \(z,z'\in\mathcal N_h\), and the exact relaxed-moment constraints \(\left|\sum_{z\in\mathcal N_h}v_{d,h,z}\mu_{dj,h}(\{z\})-m_{dj}\right|\le M_dh\) for every \(j\in J_d\). Set \(B_{d,h}(m_d)=\min\{1,V_{d,h}(m_d)+M_dh\}\), and associate each feasible vector with the canonical extension \(E_{d,h}v_{d,h}\).
\par\smallskip\noindent\emph{Inputs.} \texttt{T\_\allowbreak{}h}; \texttt{m\_\allowbreak{}d}.
\end{definition}

\begin{definition}[def:application-sensitivity-protocol]\label{def:application-sensitivity-protocol}
\par\smallskip\noindent\emph{Defined object.} \texttt{Application-\allowbreak{}facing score-\allowbreak{}law and smoothness sensitivity protocol}
\par\smallskip\noindent\emph{Construction.} Before evaluating outcomes, construct \(\mathcal G_{\mathrm{sens}}\) by taking the within-cell score laws from a complete geocoded population frame when available and otherwise including uniform-cell and prespecified population-density alternatives, then crossing those laws with separate finite grids for \(M_0\) and \(M_1\) calibrated from outcome-scale scientific bounds and negative-control spatial gradients. Report \(I(m_0,m_1)\) for every element of \(\mathcal G_{\mathrm{sens}}\), the envelope of those intervals, and a sign certificate only when every reported interval excludes zero with the same sign.
\par\smallskip\noindent\emph{Inputs.} \texttt{C\_\allowbreak{}dj}; \texttt{mu\_\allowbreak{}dj}; \texttt{M\_\allowbreak{}d}.
\end{definition}

\begin{definition}[def:verified-reconstruction-handle]\label{def:verified-reconstruction-handle}
\par\smallskip\noindent\emph{Defined object.} \texttt{Verified primal-\allowbreak{}-\allowbreak{}dual reconstruction algorithm}
\par\smallskip\noindent\emph{Construction.} \(\mathsf{VPA}\) denotes the dovetailed exact-arithmetic enumeration of conforming rational midpoint refinements, rational-nodal members of \(\mathcal F_{d,h}\) satisfying the exact continuous cell moments, unrestricted rational multipliers, rational atomic plans in \(\Pi_{\le}(\sigma)\), and outward rational interval enclosures of \(D_d\).  Given a requested rational tolerance \(\epsilon>0\), its stopping rule returns the current certificates as soon as the primal--dual gaps for both arm endpoints are at most \(\epsilon\).
\par\smallskip\noindent\emph{Inputs.} \texttt{S}; \texttt{B}; \texttt{mu\_\allowbreak{}dj}; \texttt{nu}; \texttt{M\_\allowbreak{}d}; \texttt{m\_\allowbreak{}d}; \texttt{eps}.
\end{definition}

\begin{definition}[def:confidence-box]\label{def:confidence-box}
\par\smallskip\noindent\emph{Defined object.} \texttt{Affine-\allowbreak{}hull-\allowbreak{}projected simultaneous moment region}
\par\smallskip\noindent\emph{Construction.} For each \(d\in\{0,1\}\), write \(\widehat m_d=(\widehat m_{dj})_{j\in J_d}\), define
\[
 e_{dj}=\sqrt{\frac{\log(2N/\gamma)}{2n_{dj}}},\qquad
 \rho_{d,n}=\left(\sum_{j\in J_d}e_{dj}^{2}\right)^{1/2},\qquad
 \widetilde m_d=P_{\operatorname{aff}K_d}(\widehat m_d),
\]
where \(P_{\operatorname{aff}K_d}\) is Euclidean orthogonal projection onto the affine hull of \(K_d\).  Set
\[
 \widetilde Q_{d,n}
 =\left\{q\in\operatorname{aff}K_d:
          \|q-\widetilde m_d\|_2\leq\rho_{d,n}\right\},
 \qquad
 Q_n=\widetilde Q_{0,n}\times\widetilde Q_{1,n}.
\]
\par\smallskip\noindent\emph{Inputs.} \texttt{mhat\_\allowbreak{}dj}; \texttt{n\_\allowbreak{}dj}; \texttt{gamma}; \texttt{K\_\allowbreak{}d}.
\end{definition}

\begin{definition}[def:whole-set-confidence]\label{def:whole-set-confidence}
\par\smallskip\noindent\emph{Defined object.} \texttt{Whole-\allowbreak{}set confidence inversion}
\par\smallskip\noindent\emph{Construction.} \(C_n=\bigcup_{(q_0,q_1)\in Q_n\cap(K_0\times K_1)}I(q_0,q_1)\).
\par\smallskip\noindent\emph{Inputs.} \texttt{Q\_\allowbreak{}n}.
\end{definition}

\begin{definition}[def:rectangle-envelope]\label{def:rectangle-envelope}
\par\smallskip\noindent\emph{Defined object.} \texttt{Uniform-\allowbreak{}rectangle clipping envelope}
\par\smallskip\noindent\emph{Construction.} For \(\ell_d=M_da\), define \(A_0(t)=t\), and, when \(\ell_d>0\), define
\(A_{\ell_d}(t)=[t^2-(t-\ell_d)_+^2]/(2\ell_d)\), for every \(t\in[0,1]\).
\par\smallskip\noindent\emph{Inputs.} \texttt{ell\_\allowbreak{}d}; \texttt{t}.
\end{definition}

\section{Main results}

\begin{theorem}[thm:sharp-mean-information-set]\label{thm:sharp-mean-information-set}
For every compatible \((m_0,m_1)\in K_0\times K_1\), the sharp mean-information identified set for \(\tau_B\) is \[I(m_0,m_1)=[L_1(m_1)-U_0(m_0),\ U_1(m_1)-L_0(m_0)].\] Both endpoints and every intermediate value are attained by a pair \((f_0,f_1)\in\mathcal F_0\times\mathcal F_1\), and each such pair lifts to a bounded Bernoulli potential-outcome law with the prescribed conditional means and exact observed cell means.
\end{theorem}
\begin{proof}
Put \(T(g)=\int_B g\,d\nu\) and
\[
\mathcal G_d(m)=\{g\in\mathcal F_d:A_dg=m\}.
\]
We first use the compact-fiber technique.  The domain \(S\) is compact by
\(\mathrm{ass:polygonal\mbox{-}score\mbox{-}support}\).  Every sequence in
\(\mathcal F_d\) is uniformly bounded by \(1\) and is equicontinuous with the
common modulus \(M_d\delta\), by \(\mathrm{def:surface\mbox{-}class}\).  To see
compactness without adding a hidden regularity condition, choose a countable dense
subset of \(S\), take a diagonal subsequence converging at every point of that
subset, and use a finite \(\delta\)-net together with the common modulus to show
that this subsequence is uniformly Cauchy.  Its uniform limit is continuous,
takes values in \([0,1]\), and satisfies the same Lipschitz inequality.  Hence
\(\mathcal F_d\) is compact in the uniform norm.  Moreover,
\[
 \lvert (A_dg)_j-(A_d\widetilde g)_j\rvert
 \leq \lVert g-\widetilde g\rVert_\infty,
 \qquad
 \lvert T(g)-T(\widetilde g)\rvert
 \leq \lVert g-\widetilde g\rVert_\infty,
\]
because every \(\mu_{dj}\) and \(\nu\) is a probability measure.  Thus \(A_d\)
and \(T\) are continuous.

Compatibility \(m_d\in K_d=A_d(\mathcal F_d)\), imposed by
\(\mathrm{ass:mean\mbox{-}compatibility}\), makes \(\mathcal G_d(m_d)\)
nonempty.  It is compact and convex, since it is a closed affine section of the
compact convex class \(\mathcal F_d\).  Therefore its image under \(T\) is a
compact convex subset of \(\mathbb R\), hence
\[
 T\{\mathcal G_d(m_d)\}=[L_d(m_d),U_d(m_d)],
\]
and both extrema are attained.  Since there is no cross-arm restriction, the
difference-image calculation gives
\[
\begin{split}
 I(m_0,m_1)
 &= [L_1(m_1),U_1(m_1)]-[L_0(m_0),U_0(m_0)]\\
 &= [L_1(m_1)-U_0(m_0),\ U_1(m_1)-L_0(m_0)].
\end{split}
\]
This also proves attainment of every intermediate value: every point in each
arm interval is the image of a convex combination of two feasible surfaces, and
the difference of two real intervals is the displayed interval.

It remains to distinguish this variational image from the causal identified set.
Let \(Y\in[0,1]\) be scalar, \(A\in\{0,1\}\),
\(X\in S\subset\mathbb R^2\), and let \(Z=(d,j)\) range over the finite observed
arm--cell labels.  The observable mean-information map sends a law of
\((Y,A,X,Z)\) to
\[
 \mathcal O(P)=\bigl(\{\mathcal L(X\mid Z=(d,j))=\mu_{dj}\}_{d,j},
                    \{\mathbb E[Y\mid Z=(d,j)]=m_{dj}\}_{d,j}\bigr).
\]
Define \(\Theta_I(\mathcal O(P))\) to be the values of \(\tau_B\) among all
potential-outcome laws satisfying the maintained range, Lipschitz, assignment,
and score-law restrictions and having this same observable summary.  Any such
law has conditional-mean surfaces in the two fibers, so the preceding argument
shows
\(\Theta_I(\mathcal O(P))\subseteq I(m_0,m_1)\).

For the reverse inclusion, fix any pair \((g_0,g_1)\) in the two fibers.  Choose
positive cell probabilities \(p_{dj}\) summing to one (or use the observed cell
shares when those are included in \(P\)); draw \(Z=(d,j)\) with probability
\(p_{dj}\), and draw \(X\mid Z=(d,j)\sim\mu_{dj}\).  The support, coverage,
disjointness, and boundary-null assumptions make this compatible with the sharp
side assignment \(A=d\) almost surely.  Independently draw
\(V_0,V_1\sim\mathrm{Unif}[0,1]\) conditional on \((X,Z)\), and set
\[
 Y(d)=\mathbf 1\{V_d\leq g_d(X)\},\qquad Y=Y(A).
\]
Then \(Y(d)\in[0,1]\),
\(\mathbb E[Y(d)\mid X]=g_d(X)\), and, for every observed cell,
\[
 \mathbb E[Y\mid Z=(d,j)]
 =\int_S g_d\,d\mu_{dj}=(A_dg_d)_j=m_{dj}.
\]
The causal target of this law is exactly \(\int_B(g_1-g_0)\,d\nu\).  Hence every
point of \(I(m_0,m_1)\) is generated by a bounded Bernoulli potential-outcome
law with the stipulated observed means.  This proves the reverse inclusion,
sharpness, and the asserted causal lifting.
\end{proof}
\par\smallskip\noindent \textit{Justification.} Compact moment fibers establish causal sharpness rather than a merely variational outer bound. \textit{Gap.} Qiu and Stoye impose Euclidean Lipschitz smoothness in equation \(2.1\), formulate their tight bounded isotropic extrapolation program in equation \(2.2\), and give a one-coordinate weakening in equation \(2.3\). Their framework informs the smoothness geometry here, but neither it nor Cattaneo, Titiunik, and Yu identifies a boundary functional from fixed areal integrals. \textit{Consumer.} A district-level child-stunting analysis can report a resolution-honest sign only if the complete mean-compatible interval excludes zero throughout the prespecified score-law and arm-specific-radius sensitivity grid.

\begin{theorem}[thm:clipping-transport-dual]\label{thm:clipping-transport-dual}
For every finite signed \(\sigma\) on \(S\), \[H_{M_d}(\sigma)=\min_{\pi\in\Pi_{\le}(\sigma)}\left\{\sigma^+(S)-\pi(S\times S)+M_d\int_{S\times S}\|x-y\|_2\,d\pi(x,y)\right\},\] and the minimum is attained. After the affine transformation \(\phi=2g_d-1\) and the corresponding ground-metric rescaling, this support identity is equivalent to Piccoli and Rossi's generalized \(W_1^{1,1}\) flat-metric dual. For every \(m_d\in K_d\), \[U_d(m_d)=\inf_{\lambda_d\in\mathbb R^{|J_d|}}D_d(\lambda_d;m_d),\qquad L_d(m_d)=1-U_d(\mathbf 1-m_d).\] No condition \(\sum_{j\in J_d}\lambda_{dj}=1\) is imposed. A finite minimizing \(\lambda_d\) exists whenever \(r_d(m_d)>0\). In the one-cell unit-rectangle specialization with uniform \(\mu_{dj}\), uniform boundary \(\nu\), \(M_d=1\), and \(m_d=0\), the primal value is zero but every finite multiplier has strictly positive dual objective, so multiplier attainment fails at that boundary mean.
\end{theorem}
\begin{proof}
Fix an arm \(d\), and abbreviate \(M=M_d\), \(\mathcal F=\mathcal F_d\), \(A=A_d\), \(m=m_d\), and
\[
 T(g)=\int_B g\,d\nu.
\]
The range and Lipschitz assumptions say, separately, that every \(g\in\mathcal F\) takes values in \([0,1]\) and satisfies \(\lvert g(x)-g(y)\rvert\leq M\lVert x-y\rVert_2\).  The polygonal-support assumption makes \(S\) compact and convex, hence connected.  The measures \(\mu_{dj}\) and \(\nu\) are probability measures supported on \(S\), so every integral below is finite.

We first record the compactness needed for all maxima.  Given a sequence \((g_n)\) in \(\mathcal F\), choose, for each positive integer \(k\), a finite \(k^{-1}\)-net of \(S\).  Successive subsequence extraction and a diagonal argument produce a subsequence whose values converge at every point in the union of these nets.  If \(M>0\), then, for any \(\varepsilon>0\), a sufficiently fine finite net and the common Lipschitz bound give, for all sufficiently large indices \(n,n'\),
\[
 \sup_{x\in S}\lvert g_n(x)-g_{n'}(x)\rvert
 \leq 2M\delta+\max_{z\text{ in the finite net}}
       \lvert g_n(z)-g_{n'}(z)\rvert
 <\varepsilon.
\]
If \(M=0\), connectedness makes every \(g_n\) constant and the same conclusion follows from compactness of \([0,1]\).  Thus a subsequence converges uniformly.  Its limit remains \([0,1]\)-valued and \(M\)-Lipschitz, so \(\mathcal F\) is compact in the uniform norm; it is convex by the two defining inequalities.  For every finite signed measure \(\zeta\),
\[
 \left\lvert\int_S(g-\widetilde g)\,d\zeta\right\rvert
 \leq \lVert g-\widetilde g\rVert_\infty\lvert\zeta\rvert(S),
\]
so integration is continuous.  In particular, the maximum in the support-cost definition is attained.  The same compactness also implies attainment of the arm endpoint maxima and minima whenever their exact-moment fiber is nonempty; mean compatibility supplies that nonemptiness.

We now derive the clipping-aware partial-transport formula.  Let \(\sigma=\sigma^+-\sigma^-\) be the Jordan decomposition and put
\[
 p=\sigma^+(S),\qquad q=\sigma^-(S).
\]
Suppose first that \(M>0\), define the dilation \(R(x)=2Mx\), and let
\[
 \alpha=R_\#\sigma^+,
 \qquad
 \beta=R_\#\sigma^-.
\]
For \(g\in\mathcal F\), define \(\phi(Rx)=2g(x)-1\).  Then
\[
 \lVert\phi\rVert_\infty\leq1,
 \qquad
 \lvert\phi(Rx)-\phi(Ry)\rvert
 \leq 2M\lVert x-y\rVert_2
 =\lVert Rx-Ry\rVert_2.
\]
Conversely, any function with these two bounds on \(R(S)\) gives a member of \(\mathcal F\) through \(g(x)=\{1+\phi(Rx)\}/2\).  There is no extension gap between functions on \(R(S)\) and functions on the ambient Euclidean space: the explicit formula
\[
 \Phi(z)=\inf_{u\in R(S)}\{\phi(u)+\lVert z-u\rVert_2\}
\]
is one-Lipschitz and equals \(\phi\) on \(R(S)\), while clipping \(\Phi\) to \([-1,1]\) preserves both that equality and the one-Lipschitz bound.  Therefore the bounded-Lipschitz supremum on \(R(S)\) is the ambient flat-metric supremum.

The cited Piccoli--Rossi flat-duality result, \cite[Definition 12 and Theorem 13]{PiccoliRossi2016GeneralizedWasserstein}, now gives
\[
\begin{aligned}
 H_M(\sigma)
 &=\max_{g\in\mathcal F}\int_S g\,d\sigma\\
 &=\frac{p-q}{2}
   +\frac12\sup_{\substack{\lVert\phi\rVert_\infty\leq1\\
                             \operatorname{Lip}(\phi)\leq1}}
       \int \phi\,d(\alpha-\beta)\\
 &=\frac{p-q}{2}+\frac12W_1^{1,1}(\alpha,\beta).
\end{aligned}
\]
This proves the stated affine and ground-metric equivalence, rather than merely an analogy with the generalized distance.

Apply next the cited attained submeasure representation, \cite[Definition 3 and Proposition 4]{PiccoliRossi2016GeneralizedWasserstein}.  If \(\widetilde\alpha\leq\alpha\) and \(\widetilde\beta\leq\beta\) have common mass \(t\), and \(\overline\pi\) couples them, pull \(\overline\pi\) back through \(R\times R\) to obtain \(\pi\).  Then
\[
 \pi_1\leq\sigma^+,
 \qquad
 \pi_2\leq\sigma^-,
 \qquad
 \pi(S\times S)=t,
\]
and
\[
 \int\lVert u-v\rVert_2\,d\overline\pi(u,v)
 =2M\int_{S\times S}\lVert x-y\rVert_2\,d\pi(x,y).
\]
Conversely, every \(\pi\in\Pi_{\leq}(\sigma)\) pushes forward to such an equal-mass pair of submeasures and a coupling of them.  Substituting the deletion costs \(p-t\) and \(q-t\) into the two cited identities therefore gives, equation by equation,
\[
\begin{aligned}
 H_M(\sigma)
 &=\min_{\pi\in\Pi_{\leq}(\sigma)}
   \left\{
      \frac{p-q}{2}+\frac{p+q-2\pi(S\times S)}{2}
      +M\int_{S\times S}\lVert x-y\rVert_2\,d\pi(x,y)
   \right\}\\
 &=\min_{\pi\in\Pi_{\leq}(\sigma)}
   \left\{
      p-\pi(S\times S)
      +M\int_{S\times S}\lVert x-y\rVert_2\,d\pi(x,y)
   \right\}.
\end{aligned}
\]
The minimizing submeasures in the cited representation are attained.  For those fixed submeasures, a minimizing transport coupling is also attained: a minimizing sequence has fixed finite mass on the compact set \(R(S)\times R(S)\); diagonal convergence of its integrals against a countable uniformly dense family of continuous functions gives a positive limit measure, the two marginals pass to the limit, and the continuous distance cost converges.  Pullback therefore gives an attaining \(\pi\in\Pi_{\leq}(\sigma)\).

If \(M=0\), every class member is a constant \(c\in[0,1]\), because \(S\) is connected.  Hence
\[
 H_0(\sigma)=\max_{0\leq c\leq1}c(p-q)=(p-q)_+.
\]
Every partial coupling has mass at most \(\min\{p,q\}\).  That mass is attained by the zero plan when \(pq=0\), and otherwise by
\[
 \pi=\min\{p,q\}\left(\frac{\sigma^+}{p}\right)
                    \otimes\left(\frac{\sigma^-}{q}\right).
\]
Thus the transport expression with zero ground cost equals
\[
 p-\min\{p,q\}=(p-q)_+,
\]
and its minimum is attained.  This completes the support identity for every \(M\geq0\), without any equality-of-total-mass condition on \(\sigma\).

We next dualize the exact areal moments.  Define the compact convex set
\[
 \mathcal C=\{(Ag,T(g)):g\in\mathcal F\}
 \subseteq \mathbb R^{\lvert J_d\rvert}\times\mathbb R.
\]
Compactness follows from the compactness of \(\mathcal F\) and continuity of the finitely many probability-measure integrals; convexity follows from linearity.  Let \(u=U_d(m)\).  By mean compatibility and the endpoint definition, there is \(g_0\in\mathcal F\) such that
\[
 Ag_0=m,
 \qquad
 T(g_0)=u.
\]
For any unrestricted \(\lambda\in\mathbb R^{\lvert J_d\rvert}\), the support-cost and dual-objective definitions yield
\[
\begin{aligned}
 D_d(\lambda;m)
 &=\lambda^\top m+\sup_{g\in\mathcal F}
       \left\{T(g)-\lambda^\top Ag\right\}\\
 &=\sup_{(v,t)\in\mathcal C}
       \left\{t-\lambda^\top(v-m)\right\}.
\end{aligned}
\]
Every feasible \(g\) with \(Ag=m\) consequently satisfies \(T(g)\leq D_d(\lambda;m)\).  Maximizing over the fiber and then infimizing over \(\lambda\) proves weak duality,
\[
 u\leq\inf_{\lambda\in\mathbb R^{\lvert J_d\rvert}}D_d(\lambda;m).
\]

For the reverse inequality, let \(E=\operatorname{span}(K_d-K_d)\) and work in the finite-dimensional affine Euclidean space
\[
 \mathcal L=(m+E)\times\mathbb R,
\]
which contains \(\mathcal C\).  For \(\eta>0\), the point \(y_\eta=(m,u+\eta)\) is not in \(\mathcal C\), because \(u\) is the greatest second coordinate in the fiber over \(m\).  Choose a nearest point \(z_\eta\in\mathcal C\) and write the unit vector from \(z_\eta\) to \(y_\eta\) as
\[
 n_\eta=(a_\eta,b_\eta)\in E\times\mathbb R.
\]
For every \(c\in\mathcal C\), comparison of \(z_\eta\) with \(z_\eta+s(c-z_\eta)\), followed by differentiation of squared distance at \(s=0\), gives
\[
 n_\eta^\top(c-z_\eta)\leq0.
\]
Since \(n_\eta^\top(y_\eta-z_\eta)>0\), comparison with \((m,u)\in\mathcal C\) gives
\[
 b_\eta\eta=n_\eta^\top\{y_\eta-(m,u)\}>0,
\]
so \(b_\eta>0\).  For \(c=(v,t)\in\mathcal C\), the same separating inequality implies
\[
 a_\eta^\top(v-m)+b_\eta(t-u-\eta)\leq0.
\]
Set \(\lambda_\eta=-a_\eta/b_\eta\), viewed as an unrestricted vector in \(\mathbb R^{\lvert J_d\rvert}\).  Then
\[
 t-\lambda_\eta^\top(v-m)\leq u+\eta
 \quad\text{for every }(v,t)\in\mathcal C,
\]
and hence
\[
 D_d(\lambda_\eta;m)\leq u+\eta.
\]
Letting \(\eta\downarrow0\) and combining this inequality with weak duality proves
\[
 U_d(m)=\inf_{\lambda\in\mathbb R^{\lvert J_d\rvert}}D_d(\lambda;m).
\]
The separating vector ranged over the whole multiplier space; neither the argument nor the dual-objective definition imposes \(\sum_{j\in J_d}\lambda_{dj}=1\).

Suppose now that \(r_d(m)>0\).  By the relative-interior-radius definition, there is a \(\rho>0\) such that
\[
 m+\rho\mathbb B_E\subseteq K_d.
\]
Apply the preceding normalized separation with \(\eta=1/n\).  A subsequence of the unit normals \((a_n,b_n)\) converges to some unit vector \((a,b)\in E\times\mathbb R\).  Passing to the limit in the separating inequality gives
\[
 a^\top(v-m)+b(t-u)\leq0
 \quad\text{for every }(v,t)\in\mathcal C.
\]
The limit cannot have \(b=0\).  Indeed, then \(\lVert a\rVert_2=1\), while \(m+\rho a\in K_d\); choosing \(g\in\mathcal F\) with \(Ag=m+\rho a\) would give \(\rho=a^\top(Ag-m)\leq0\), a contradiction.  Therefore \(b>0\).  With \(\lambda=-a/b\),
\[
 t-\lambda^\top(v-m)\leq u
 \quad\text{for every }(v,t)\in\mathcal C,
\]
so \(D_d(\lambda;m)\leq u\).  Weak duality gives the reverse inequality, and hence this finite \(\lambda\) attains the dual minimum.

For the lower endpoint, probability normalization of every \(\mu_{dj}\) and of \(\nu\) makes \(g\mapsto1-g\) a bijection of \(\mathcal F\) satisfying
\[
 A(1-g)=\mathbf 1-Ag,
 \qquad
 T(1-g)=1-T(g).
\]
In particular, \(m\in K_d\) implies \(\mathbf1-m\in K_d\), and reflection of the exact-moment fiber yields
\[
 L_d(m)=1-U_d(\mathbf1-m).
\]

Finally, specialize to the one-cell rectangle
\[
 S=C_{dj}=[0,1]\times[0,1],
 \qquad
 B=\{0\}\times[0,1],
\]
with uniform area law \(\mu_{dj}\), uniform length law \(\nu\), \(M=1\), and \(m=0\).  If \(g\) is feasible, then \(g\geq0\) and \(\int_Sg\,d\mu_{dj}=0\).  Were \(g\) positive at any point, continuity would make it positive on a relative neighborhood of positive uniform area, contradicting the zero integral.  Thus \(g\equiv0\), including on \(B\), and
\[
 U_d(0)=0.
\]
For any finite scalar multiplier \(\lambda\) and any \(t\in(0,1]\), the cone
\[
 g_t(x,y)=(t-x)_+
\]
is \([0,1]\)-valued and one-Lipschitz, with
\[
 \int_Bg_t\,d\nu=t,
 \qquad
 \int_Sg_t\,d\mu_{dj}
 =\int_0^t(t-x)\,dx
 =\frac{t^2}{2}.
\]
The support definition therefore implies
\[
 D_d(\lambda;0)=H_1(\nu-\lambda\mu_{dj})
 \geq t-\frac{\lambda t^2}{2}.
\]
If \(\lambda\leq0\), take any \(t\in(0,1]\); if \(\lambda>0\), take \(0<t<\min\{1,2/\lambda\}\).  In either case the right-hand side is strictly positive.  Thus every finite multiplier has \(D_d(\lambda;0)>0\), although the proved value identity gives \(\inf_\lambda D_d(\lambda;0)=U_d(0)=0\).  No finite multiplier attains the infimum at this boundary mean, whereas the primal surface and the partial-transport minimizer do attain their respective optima.
\end{proof}
\par\smallskip\noindent \textit{Justification.} The theorem combines an attained clipping-aware partial-transport support cost with an exact unrestricted-multiplier representation and identifies relative-interior attainment as the precise positive computational regime. \textit{Gap.} Piccoli and Rossi's Definition 12 and Theorem 13 supply flat duality, and their Definition 3 and Proposition 4 supply attained submeasure transport, but neither result incorporates exact areal side moments or their unrestricted finite-dimensional multipliers. Qiu and Stoye's equations \(2.1\)--\(2.3\) inform the bounded isotropic geometry but concern a different boundary-to-policy extrapolation problem. \textit{Consumer.} Certified endpoint software may solve an attained partial-transport problem at each multiplier, use finite minimizers for relative-interior means, and recognize the zero-mean rectangle as genuine multiplier nonattainment rather than numerical failure.

\begin{theorem}[thm:closed-form-specializations]\label{thm:closed-form-specializations}
Three complete planar reductions hold. First, for one uniform cell \([0,a]\times[0,1]\) and uniform boundary law on \(\{0\}\times[0,1]\), \(U_d(m_d)=\max\{t\in[0,1]:A_{\ell_d}(t)\le m_d\}\), equivalently \[U_d(m_d)=\begin{cases}1,&m_d\ge A_{\ell_d}(1),\\ \sqrt{2\ell_dm_d},&m_d<A_{\ell_d}(1)\ \text{and }m_d\le\ell_d/2,\\ m_d+\ell_d/2,&\text{otherwise},\end{cases}\qquad L_d(m_d)=1-U_d(1-m_d).\] Second, on the unit rectangle with \(M_d=1\) and horizontal density \(\rho_\star\), \[U_d(m_d)=\begin{cases}\sqrt{4m_d/3},&0\le m_d\le3/16,\\ \sqrt{3/2+4m_d}-1,&3/16\le m_d\le5/8,\\1,&5/8\le m_d\le1,\end{cases}\qquad L_d(m_d)=1-U_d(1-m_d).\] Third, on the triangle with vertices \((-1,0),(1,0),(0,1)\), split at \(x=0\), with uniform area laws and boundary density \(w(0,y)=2(1-y)\), one has \([L_d(m_d),U_d(m_d)]=[m_d-M_d/3,m_d+M_d/3]\) whenever \(0\le M_d\le3/4\) and \(2M_d/3\le m_d\le1-2M_d/3\).
\end{theorem}
\begin{proof}
We use one-dimensional symmetrization for the two rectangular cases.  Existence
of the stated two-dimensional uniform probability law entails \(a>0\), which
justifies division by \(a\).  In the uniform rectangle let
\[
 \bar g(x)=\int_0^1g(x,y)\,dy.
\]
Then \(0\leq\bar g\leq1\),
\(\lvert\bar g(x)-\bar g(x')\rvert\leq M_d\lvert x-x'\rvert\), its uniform
average on \([0,a]\) equals the cell mean, and
\(\bar g(0)=\int_Bg\,d\nu\).  Conversely, a one-dimensional function with
these properties gives the feasible surface \(g(x,y)=\bar g(x)\) on the cell,
and the clipped McShane formula extends it to the rest of \(S\) without changing
its values there.  Thus this reduction is exact, not an outer bound.

Fix a boundary height \(t\).  The Lipschitz and range restrictions imply
\[
 \bar g(x)\geq(t-M_dx)_+.
\]
Equality is feasible, and hence, for \(\ell_d=M_da>0\), the least possible cell
mean is
\[
\begin{split}
 \frac1a\int_0^a(t-M_dx)_+\,dx
 &=\frac{t^2-(t-M_da)_+^2}{2M_da}\\
 &=A_{\ell_d}(t).
\end{split}
\]
When \(\ell_d=0\), the function is constant and the least mean is \(t=A_0(t)\),
which explains the required zero-radius completion of the envelope definition.
Reflection around \(1/2\) shows that the greatest mean at height \(t\) is
\(1-A_{\ell_d}(1-t)\).  Convex combinations of the two extremizers keep their
boundary height and fill the interval between these mean values.  At the largest
feasible \(t<1\) the lower inequality binds, while \(t=1\) is feasible exactly
when \(m_d\geq A_{\ell_d}(1)\).  Therefore
\[
 U_d(m_d)=\max\{t\in[0,1]:A_{\ell_d}(t)\leq m_d\}.
\]
For \(\ell_d>0\), the envelope has the two branches
\[
 A_{\ell_d}(t)=
 \begin{cases}
 t^2/(2\ell_d),&t\leq\ell_d,\\
 t-\ell_d/2,&t\geq\ell_d.
 \end{cases}
\]
Inverting them, including the clipping point \(t=1\), gives
\[
U_d(m_d)=\begin{cases}
1,&m_d\geq A_{\ell_d}(1),\\
\sqrt{2\ell_dm_d},&m_d<A_{\ell_d}(1)\ \text{and }m_d\leq\ell_d/2,\\
m_d+\ell_d/2,&\text{otherwise}.
\end{cases}
\]
The same display remains correct at \(\ell_d=0\) under the corrected definition,
and reflection gives \(L_d(m_d)=1-U_d(1-m_d)\).

For the nonuniform rectangle the same reduction applies, now with horizontal
density \(\rho_\star\).  At height \(t\), direct integration of the extremal cone
\((t-x)_+\) gives
\[
 A_\star(t)=
 \begin{cases}
 \displaystyle\frac32\int_0^t(t-x)\,dx=\frac{3t^2}{4},&0\leq t\leq\frac12,\\[6pt]
 \displaystyle\frac32\int_0^{1/2}(t-x)\,dx
 +\frac12\int_{1/2}^{t}(t-x)\,dx
 =\frac{t^2}{4}+\frac t2-\frac18,&\frac12\leq t\leq1.
 \end{cases}
\]
The transition means are \(A_\star(1/2)=3/16\) and
\(A_\star(1)=5/8\).  Solving \(A_\star(t)=m_d\) on the two branches and clipping
at one proves
\[
U_d(m_d)=\begin{cases}
\sqrt{4m_d/3},&0\leq m_d\leq3/16,\\
\sqrt{3/2+4m_d}-1,&3/16\leq m_d\leq5/8,\\
1,&5/8\leq m_d\leq1,
\end{cases}
\qquad L_d(m_d)=1-U_d(1-m_d).
\]

Finally use an explicit coupling for the triangle.  On the right-hand triangular
cell \(T=\{(x,y):x\geq0,\ x+y\leq1\}\), the uniform area density is \(2\).  The
projection \(P(x,y)=(0,y)\) has pushforward density
\[
 \int_0^{1-y}2\,dx=2(1-y)=w(0,y),
\]
so \(P_\#\mu_T=\nu\).  Since the centroid of \(T\) has first coordinate \(1/3\),
the Lipschitz inequality yields
\[
 \left|\int_Bg\,d\nu-\int_Tg\,d\mu_T\right|
 =\left|\int_T\{g(P(x,y))-g(x,y)\}\,d\mu_T\right|
 \leq M_d\int_Tx\,d\mu_T=\frac{M_d}{3}.
\]
The affine functions
\[
 g_+(x,y)=m_d+\frac{M_d}{3}-M_dx,
 \qquad
 g_-(x,y)=m_d-\frac{M_d}{3}+M_dx
\]
have cell mean \(m_d\) and boundary means \(m_d+M_d/3\) and
\(m_d-M_d/3\), respectively.  On \(T\), both remain in \([0,1]\) whenever
\(2M_d/3\leq m_d\leq1-2M_d/3\); this interval is nonempty when
\(0\leq M_d\leq3/4\).  Clipping their affine continuation to \([0,1]\) on all
of \(S\) preserves the Lipschitz constant and does not alter their values on
\(T\).  They attain the two coupling bounds, proving
\[
 [L_d(m_d),U_d(m_d)]=[m_d-M_d/3,m_d+M_d/3].
\]
Reflecting the horizontal coordinate proves the same formula for the left-hand
cell.  This completes all three exact reductions.
\end{proof}
\par\smallskip\noindent \textit{Justification.} Closed forms expose clipping transitions, density sensitivity, and nonconstant boundary weights that an abstract optimization display conceals. \textit{Gap.} The closest boundary papers provide estimators or a coordinate projection radius, while the flat-metric literature does not invert exact areal means into these endpoint laws. \textit{Consumer.} The formulas give immediate fixed-resolution sign and informativeness checks before any general-purpose numerical solve.

\begin{theorem}[thm:relative-interior-programs]\label{thm:relative-interior-programs}
For the fixed nested midpoint refinement and the lexicographically tie-broken projections, \(\|x-p_{dj,h}(x)\|_2\le h\) on every \(C_{dj}\) and \(\|x-p_{B,h}(x)\|_2\le h\) on \(B\). Consequently every \(g_d\in\mathcal F_d\) yields a feasible nodal vector \(v_{d,h,z}=g_d(z)\), its cell and boundary pushforward errors are at most \(M_dh\), and \(U_d(m_d)\le B_{d,h}(m_d)\) for every \(m_d\in K_d\). For every compatible mean, \(B_{d,h}(m_d)\to U_d(m_d)\) as \(h\downarrow0\). For every atomic optimizer \(v_{d,h}\), the canonical extension satisfies \[\left\|A_d(E_{d,h}v_{d,h})-m_d\right\|_2\le2M_dh\sqrt{|J_d|},\qquad \left|\int_BE_{d,h}v_{d,h}\,d\nu-\sum_{z\in\mathcal N_h}v_{d,h,z}\nu_h(\{z\})\right|\le M_dh.\] If \(r_d(m_d)>0\), relative-interior endpoint stability gives \(U_d(A_dE_{d,h}v_{d,h})\le U_d(m_d)+\|A_dE_{d,h}v_{d,h}-m_d\|_2/r_d(m_d)\); adding the boundary projection error and the \(M_dh\) correction in \(B_{d,h}\) yields exactly \[0\le B_{d,h}(m_d)-U_d(m_d)\le2M_dh\left(1+\frac{\sqrt{|J_d|}}{r_d(m_d)}\right).\] The exact-moment piecewise-affine fiber is then nonempty for all sufficiently small \(h\), \(P_{d,h}(m_d)\le U_d(m_d)\), and \[0\le U_d(m_d)-P_{d,h}(m_d)\le\varepsilon_{d,h}\left(1+\frac{\sqrt{|J_d|}}{r_d(m_d)}\right),\qquad \varepsilon_{d,h}=O(\sqrt h).\] The reflected programs give the corresponding lower-endpoint bounds and shrinking a posteriori effect-set gaps when certified finite-program optimization errors vanish.
\end{theorem}
\begin{proof}
Fix an arm \(d\), and suppress the arm subscript when no ambiguity can
result.  Write \(J=|J_d|\), \(M=M_d\), \(A=A_d\), \(m=m_d\),
\(U=U_d\), and \(T(g)=\int_B g\,d\nu\).  All norms on the moment
space below are Euclidean.  The proof has four parts: a geometric projection
argument, a compactness argument, a support-function moment correction, and a
clipping-preserving finite-element approximation.

First consider the projections.  Because the initial triangulation conforms to
every cell boundary and to \(B\), so does every midpoint refinement.  If
\(x\in C_{dj}\), some triangle lying in \(C_{dj}\) and containing \(x\) has a
vertex in \(\mathcal N_h\cap C_{dj}\) at distance at most its diameter.  The
same argument on a boundary edge applies to \(x\in B\).  Nearest-node
selection, with lexicographic tie breaking, is Borel measurable and therefore
satisfies
\[
 \|x-p_{dj,h}(x)\|_2\le h,
 \qquad
 \|x-p_{B,h}(x)\|_2\le h.                                      \tag{1}
\]
Here and below the support condition on \(\mu_{dj}\) is used when integrating
the first inequality.  For any \(g\in\mathcal F_d\), set
\(v_z=g(z)\) at every \(z\in\mathcal N_h\).  The range restriction and the
\(M\)-Lipschitz restriction give all nodal constraints.  The pushforward
definitions and (1) give, for every \(j\in J_d\),
\begin{align}
 \left|\sum_z v_z\mu_{dj,h}(\{z\})-m_{dj}\right|
 &=\left|\int_{C_{dj}}\{g(p_{dj,h}(x))-g(x)\}\,d\mu_{dj}(x)\right|
 \le Mh,                                                       \tag{2}\\
 \left|\sum_z v_z\nu_h(\{z\})-T(g)\right|
 &=\left|\int_B\{g(p_{B,h}(x))-g(x)\}\,d\nu(x)\right|
 \le Mh.                                                       \tag{3}
\end{align}
Both laws in these displays have total mass one.  Compatibility makes the
exact moment fiber nonempty.  It is uniformly closed inside the uniformly
bounded, common-modulus class \(\mathcal F_d\), and that class is compact in
the uniform norm: to see this without importing a compactness theorem, choose
successive finite \(1/k\)-nets of compact \(S\), diagonalize a bounded sequence
on the countable union of the nets, and use the common Lipschitz modulus to
upgrade pointwise convergence on that dense set to uniform convergence.  The
uniform limit remains in \(\mathcal F_d\).  Thus an optimizer \(g^*\) attaining
\(U(m)\) exists.  Applying (2)--(3) to \(g^*\) proves
\[
 U(m)\le V_{d,h}(m)+Mh.
\]
Since \(0\le U(m)\le1\), the clipping in the definition of \(B_{d,h}\)
preserves the inequality:
\[
 U(m)\le B_{d,h}(m).                                         \tag{4}
\]

Conversely, let \(v\) be feasible for the atomic program.  Pairwise nodal
feasibility implies
\[
 \min_{z'\in\mathcal N_h}\{v_{z'}+M\|z-z'\|_2\}=v_z
 \quad(z\in\mathcal N_h).
\]
Indeed the left side is at most \(v_z\) by choosing \(z'=z\), and every other
term is at least \(v_z\) by the pairwise constraint.  The minimum of
\(M\)-Lipschitz functions is \(M\)-Lipschitz, and clipping by
\(t\mapsto\min\{1,\max\{0,t\}\}\) is a contraction.  Hence
\(e_h=E_{d,h}v\) belongs to \(\mathcal F_d\) and equals \(v\) at every node.
Combining the relaxed moment constraint with (1) yields
\[
 \|Ae_h-m\|_2\le 2Mh\sqrt J,                                 \tag{5}
\]
while the boundary calculation yields
\[
 \left|T(e_h)-\sum_zv_z\nu_h(\{z\})\right|\le Mh.            \tag{6}
\]
The finite atomic feasible set is a nonempty compact polytope, so an atomic
optimizer exists.  Along an arbitrary sequence \(h_n\downarrow0\), take such
optimizers and their extensions.  The elementary compactness argument above
gives a uniformly convergent subsequence with limit \(e\in\mathcal F_d\).
Equations (5)--(6) imply \(Ae=m\) and that the corresponding atomic objectives
converge to \(T(e)\le U(m)\).  Since every subsequence has a further subsequence
with this upper-limit property,
\(\limsup_{h\downarrow0}V_{d,h}(m)\le U(m)\).  Moreover
\(0\le B_{d,h}(m)-V_{d,h}(m)\le Mh\).  Together with (4), this proves
\[
 B_{d,h}(m)\longrightarrow U(m)                              \tag{7}
\]
at every compatible mean, including a mean on the relative boundary of the
moment body.  This step is a value-side compactness argument and makes no exact
finite-mesh feasibility claim.

We next derive the relative-interior quantitative bound.  Put
\(r=r_d(m)>0\) and \(q=Ae_h\).  Since \(q,m\in K_d\), the vector
\(q-m\) belongs to \(E_d\).  If \(q\ne m\), put
\(a=(q-m)/\|q-m\|_2\).  The definition of \(r\) gives
\(m-ra\in K_d\), and the identity
\[
 m=\frac{r}{r+\|q-m\|_2}q
   +\frac{\|q-m\|_2}{r+\|q-m\|_2}(m-ra)                     \tag{8}
\]
holds.  The endpoint map \(U\) is concave on \(K_d\): convex combinations of
two feasible surfaces have the corresponding convex-combination moments and
target values.  Applying concavity to (8) and using \(0\le U\le1\) gives
\[
 U(q)\le U(m)+\frac{\|q-m\|_2}{r}.                           \tag{9}
\]
The same inequality is immediate if \(q=m\).  For an atomic optimizer,
(6), feasibility of \(e_h\) at moment \(q\), and (9) imply
\[
 V_{d,h}(m)\le T(e_h)+Mh
 \le U(m)+\frac{\|q-m\|_2}{r}+Mh.
\]
Adding the final \(Mh\) in the definition of \(B_{d,h}\), noting that clipping
at one can only decrease the result, and using (5), proves the stated exact
outer error:
\[
 0\le B_{d,h}(m)-U(m)
 \le 2Mh\left(1+\frac{\sqrt J}{r}\right).                   \tag{10}
\]

For the inner correction, set
\[
 K_{d,h}=A(\mathcal F_{d,h}),
 \qquad
 \delta_h=\sqrt J\,\varepsilon_{d,h}.
\]
The nodal feasible set defining \(\mathcal F_{d,h}\) is compact and convex.
Vertex range bounds extend to each triangle by barycentricity.  On any line
segment in convex \(S\), integration of the piecewise-constant gradient and
the trianglewise gradient bound give the global Euclidean Lipschitz inequality.
Thus
\[
 \mathcal F_{d,h}\subseteq\mathcal F_d,
 \qquad K_{d,h}\subseteq K_d,                               \tag{11}
\]
and \(K_{d,h}\) is compact and convex.  Let \(s_C(a)=\sup_{x\in C}a^\top x\)
denote a support function.  For \(a\in E_d\), choose a support maximizer in
\(\mathcal F_d\) and a best uniform approximant in \(\mathcal F_{d,h}\).
Because every \(\mu_{dj}\) is a probability law,
\(\|A(g-q)\|_2\le\sqrt J\|g-q\|_\infty\).  Therefore
\begin{align}
 s_{K_{d,h}}(a)
 &\ge s_{K_d}(a)-\delta_h\|a\|_2\\
 &\ge a^\top m+(r-\delta_h)\|a\|_2.                        \tag{12}
\end{align}
The second inequality uses
\(m+r\mathbb B_{E_d}\subseteq K_d\).  Equation (12) implies
\[
 m+(r-\delta_h)\mathbb B_{E_d}\subseteq K_{d,h}             \tag{13}
\]
whenever \(\delta_h<r\).  For completeness, if a point of the ball in (13)
were outside the compact convex set \(K_{d,h}\), its nearest point in the
finite-dimensional affine space \(m+E_d\) would furnish a separating vector
\(a\in E_d\), contradicting (12).  This is the support-function separation
step.

Let \(g^*\) attain \(U(m)\), and choose
\(g_h\in\mathcal F_{d,h}\) with
\(\|g_h-g^*\|_\infty\le\varepsilon_{d,h}\).  Put
\(e=Ag_h-m\); then \(e\in E_d\) and \(\|e\|_2\le\delta_h\).  If
\(\delta_h=0\), the surface \(g_h\) is exactly moment feasible.  Otherwise set
\(t=\delta_h/r\).  Since
\[
 \left\|-\frac{1-t}{t}e\right\|_2
 \le r-\delta_h,
\]
(13) supplies \(q_h\in\mathcal F_{d,h}\) such that
\[
 Aq_h=m-\frac{1-t}{t}e.
\]
The convex combination \(f_h=(1-t)g_h+tq_h\) then satisfies
\(Af_h=m\).  Positivity and normalization of \(\nu\), together with the range
restriction, imply \(0\le T(q_h),T(g^*)\le1\).  Hence
\begin{align}
 T(f_h)
 &=(1-t)T(g_h)+tT(q_h)\\
 &\ge(1-t)\{U(m)-\varepsilon_{d,h}\}\\
 &\ge U(m)-\varepsilon_{d,h}-\frac{\delta_h}{r}.            \tag{14}
\end{align}
Since \(\delta_h<r\) for all sufficiently small \(h\), this proves eventual
nonemptiness of the exact piecewise-affine fiber.  Inclusion (11) gives
\(P_{d,h}(m)\le U(m)\), while (14) gives
\[
 0\le U(m)-P_{d,h}(m)
 \le\varepsilon_{d,h}\left(1+\frac{\sqrt J}{r}\right).     \tag{15}
\]

It remains to prove, rather than assume, the claimed approximation rate.  All
children of a triangle in a midpoint refinement are similar to their parent.
Consequently, for the edge-row matrix \(R_T\) of a triangle \(T\),
\[
 C_0:=\sup_{h,T\in\mathcal T_h}
       \frac{\operatorname{diam}(T)\|R_T^{-1}\|_{\mathrm{op}}}{\sqrt2}<\infty.
                                                                    \tag{16}
\]
Suppose first that \(M>0\).  The clipped McShane formula extends
\(g\in\mathcal F_d\) to an \([0,1]\)-valued \(M\)-Lipschitz function
\(\widetilde g\) on \(\mathbb R^2\).  For a standard Gaussian vector
\(Z\in\mathbb R^2\), define
\[
 u_\delta(x)=\mathbb E\{\widetilde g(x+\delta Z)\}.
\]
Direct coupling and the polar density
\(s e^{-s^2/2}\mathbf 1\{s\ge0\}\) give
\[
 \|u_\delta-g\|_\infty
 \le M\delta\,\mathbb E\|Z\|_2
 =M\delta\sqrt{\frac\pi2}.                                \tag{17}
\]
The function \(\widetilde g\) has distributional first derivatives bounded in
Euclidean norm by \(M\): this follows directly by taking bounded directional
difference quotients, testing against a smooth compactly supported function,
and passing to a weak limit in the integration-by-parts identity.  Convolution
with the Gaussian kernel is smooth.  Differentiating the convolution of that
bounded weak gradient in a unit direction \(a\) gives
\[
 D_a\nabla u_\delta(x)
 =\frac1\delta\mathbb E\!\left[
       \nabla\widetilde g(x+\delta Z)(Z^\top a)\right].
\]
Since \(\mathbb E|Z^\top a|=\sqrt{2/\pi}\),
\[
 \|D^2u_\delta\|_{\mathrm{op}}
 \le\frac{M}{\delta}\sqrt{\frac2\pi}.                     \tag{18}
\]

Let \(I_hu_\delta\) be the nodal affine interpolant.  Taylor expansion from one
vertex along the other two edges shows that the two-component remainder vector
has norm at most
\(\|D^2u_\delta\|_{\mathrm{op}}\operatorname{diam}(T)^2/\sqrt2\).
Multiplication by \(R_T^{-1}\), followed by (16)--(18), yields
\[
 \|\nabla(I_hu_\delta)|_T\|_2\le M(1+\vartheta),
 \qquad
 \vartheta=C_0\sqrt{\frac2\pi}\frac h\delta.              \tag{19}
\]
Barycentric interpolation preserves the range and, using only that
\(u_\delta\) is \(M\)-Lipschitz, also gives
\[
 \|I_hu_\delta-u_\delta\|_\infty\le Mh.                    \tag{20}
\]
Contract the interpolant about \(1/2\):
\[
 q_h=\frac12+\frac{I_hu_\delta-1/2}{1+\vartheta}.
\]
Equations (19)--(20) show that \(q_h\in\mathcal F_{d,h}\).  Since
\(|I_hu_\delta-1/2|\le1/2\), equations (17)--(20) imply
\[
 \|q_h-g\|_\infty
 \le M\delta\sqrt{\frac\pi2}+Mh+\frac{\vartheta}{2}.       \tag{21}
\]
The right side is uniform over \(g\in\mathcal F_d\).  Taking
\(\delta=\sqrt h\) in (21) proves
\[
 \varepsilon_{d,h}=O(\sqrt h),                              \tag{22}
\]
where the constant depends only on \(M_d\) and the fixed initial mesh.  If
\(M=0\), convexity of \(S\) makes every class member constant, and constants
belong to every \(\mathcal F_{d,h}\); hence
\(\varepsilon_{d,h}=0\).  This completes the inner-program proof.

Finally, \(g\mapsto1-g\) bijects the fiber at \(m_d\) with the fiber at
\(\mathbf1-m_d\), preserves every approximation class, and gives
\(L_d(m_d)=1-U_d(\mathbf1-m_d)\).  Applying (4), (7), (10), and (15) to the
reflected programs supplies the lower-endpoint bounds.  Combining the two
armwise endpoint brackets by interval subtraction gives an outer bracket and
an inner attainable interval for \(I(m_0,m_1)\); their endpoint gaps are sums
of the corresponding armwise gaps and therefore vanish when the certified
finite-program optimization errors vanish.  This proves every assertion of the
theorem.
\end{proof}
\par\smallskip\noindent \textit{Justification.} The result separates universal value-side approximation from exact primal recovery and quantifies precisely where relative-interior slack repairs continuous moments. \textit{Gap.} Mula and Nouy prove moment-hierarchy convergence, and Baldi and Mourrain quantify pseudo-moment approximation, but neither supplies this exact response-surface moment correction or its clipping-preserving endpoint gap. \textit{Consumer.} A practitioner receives a monotonically valid outer interval at every compatible mean and a convergent exact-feasibility inner certificate away from exposed faces.

\begin{theorem}[thm:confidence-box-inversion]\label{thm:confidence-box-inversion}
Under conditional cell-sample independence,
\[
\inf\Pr\{I(m_0,m_1)\subseteq C_n\}\geq1-\gamma,
\]
where the infimum ranges over all maintained bounded cell-outcome laws with the
fixed score laws and compatible means. Whenever
\(Q_n\cap(K_0\times K_1)\neq\varnothing\), \(C_n\) is an attained compact
interval equal to the effect image of all response-surface pairs whose joint
moment vector lies in \(Q_n\). If \(r_d(m_d)>0\) and
\(E_d=\mathbb R^{|J_d|}\) for both arms, and
\(\min_{d,j}n_{dj}\to\infty\), then the Hausdorff distance between \(C_n\) and
\(I(m_0,m_1)\) converges to zero in probability.
\end{theorem}
\begin{proof}
Fix any maintained data-generating law with compatible means \(m_d\in K_d\), \(d\in\{0,1\}\), and conditional cell-sample independence.  By \(\mathrm{thm:affine\mbox{-}hull\mbox{-}projected\mbox{-}confidence\mbox{-}inversion}\), applied to the projected product region \(Q_n=\widetilde Q_{0,n}\times\widetilde Q_{1,n}\) from \(\mathrm{def:confidence\mbox{-}box}\) and its inversion \(C_n\) from \(\mathrm{def:whole\mbox{-}set\mbox{-}confidence}\),
\[
 \Pr\!\left\{I(m_0,m_1)\subseteq C_n\right\}\ge 1-\gamma.
\]
The cited theorem establishes this inequality uniformly over all maintained \([0,1]\)-valued cell-outcome laws with the fixed score laws and compatible means.  Taking the infimum over that same class therefore gives
\[
 \inf_{\mathbb P}\Pr_{\mathbb P}\!\left\{I(m_0,m_1)\subseteq C_n\right\}\ge 1-\gamma.
\]

The same theorem states that, on the event
\[
 Q_n\cap(K_0\times K_1)\ne\varnothing,
\]
the set \(C_n\) is an attained compact interval and equals the image of all pairs \((g_0,g_1)\in\mathcal F_0\times\mathcal F_1\) satisfying \((A_0g_0,A_1g_1)\in Q_n\), under the continuous linear effect map
\[
 (g_0,g_1)\longmapsto \int_B(g_1-g_0)\,d\nu.
\]
This is exactly the target's attained-interval and exact-effect-image assertion.

Finally, suppose that \(r_d(m_d)>0\) and \(E_d=\mathbb R^{|J_d|}\) for both arms, and that \(\min_{d,j}n_{dj}\to\infty\).  The relative-interior consistency part of \(\mathrm{thm:affine\mbox{-}hull\mbox{-}projected\mbox{-}confidence\mbox{-}inversion}\) requires only
\[
 r_d(m_d)>0\quad(d=0,1),
 \qquad
 \min_{d,j}n_{dj}\longrightarrow\infty,
\]
and imposes no condition that \(E_d\) equal the ambient space.  Hence it applies under the target's stronger hypotheses and yields
\[
 \Pr\!\left\{Q_n\cap(K_0\times K_1)\ne\varnothing\right\}\longrightarrow1,
 \qquad
 d_{\mathrm H}\!\left(C_n,I(m_0,m_1)\right)\xrightarrow{\mathbb P}0.
\]
The second convergence is the target's final assertion; the first also ensures that the possibly empty finite-sample definition of \(C_n\) is nonempty with probability tending to one.  Thus all three claims follow.
\end{proof}
\par\smallskip\noindent \textit{Justification.} Affine-hull projection makes simultaneous whole-set inference valid for lower-dimensional moment bodies, covers both endpoints and every value in the sharp set, and requires no differentiation of a potentially nonsmooth endpoint map. \textit{Gap.} Chernozhukov, Hong, and Tamer give asymptotic criterion inversion, while Qiu and Stoye use an exact-normal approximation and Cattaneo, Titiunik, and Yu infer a point-valued boundary curve from unit-level data; none supplies finite-sample inversion of fixed areal means in the relative affine geometry of the compatible moment body. \textit{Consumer.} A sign certificate remains valid after simultaneous sampling uncertainty in every district-level stunting mean is projected into the compatible affine geometry and propagated through the entire sharp set.

\begin{proposition}[prop:sanity-reductions]\label{prop:sanity-reductions}
The kernel obeys three known-case reductions. If \(M_d=0\), compatibility forces every coordinate of \(m_d\) to equal a common constant and \(L_d(m_d)=U_d(m_d)\) equals that constant. In the one-cell uniform rectangle with \(a=M_d=1\) and \(m_d=1/2\), each arm interval is \([0,1]\), so equal arm means give \(I(m_0,m_1)=[-1,1]\). If \(m_1=7/8\) and \(m_0=1/32\) in that same geometry, then \(I(m_0,m_1)=[1/4,1]\).
\end{proposition}
\begin{proof}
If \(M_d=0\), every member of \(\mathcal F_d\) is a constant \(c\in[0,1]\).
Because every \(\mu_{dj}\) is a probability measure, compatibility is equivalent
to \(m_{dj}=c\) for every \(j\), and the boundary integral is also \(c\).  Hence
\(L_d(m_d)=U_d(m_d)=c\).

For the unit uniform rectangle, \(\ell_d=1\) and
\(A_1(1)=1/2\).  At \(m_d=1/2\), the closed form gives \(U_d=1\), while reflection
gives \(L_d=1-U_d(1/2)=0\).  Thus each arm interval is \([0,1]\), and the exact
difference interval is \([-1,1]\).

At \(m_1=7/8\), the same formula gives \(U_1=1\) and
\[
 L_1=1-U_1(1/8)=1-\sqrt{2(1/8)}=\frac12.
\]
At \(m_0=1/32\), it gives
\[
 U_0=\sqrt{2(1/32)}=\frac14,
 \qquad L_0=1-U_0(31/32)=0.
\]
Substitution into the sharp effect-set formula yields
\[
 I(m_0,m_1)=[L_1-U_0,U_1-L_0]=[1/4,1].
\]
All three reductions therefore follow from the proved endpoint formulas and the
sharp-set theorem.
\end{proof}
\par\smallskip\noindent \textit{Justification.} The reductions check the zero-smoothness limit, the logically vacuous fixed-resolution case, and a nonvacuous sign-identifying case. \textit{Gap.} These are required internal reductions rather than novelty claims; no closest-paper improvement is asserted. \textit{Consumer.} They serve as analytic regression tests for any implementation of the general endpoint program.

\begin{theorem}[thm:terminating-verified-reconstruction]\label{thm:terminating-verified-reconstruction}
Let \(S\) be a compact convex rational polygon, let the cells and \(B\) have
rational polygonal geometry, let every \(\mu_{dj}\) have a rational
piecewise-constant density, and let \(d\nu=w\,d\mathcal H^1\), where \(w\) is
nonnegative, normalized, and rational piecewise affine on \(B\).  If \(M_d\geq0\),
\(m_d\in K_d\), and \(\epsilon>0\) are rational and \(r_d(m_d)>0\), then there is
a terminating exact-arithmetic implementation of \(\mathsf{VPA}\) which returns
rational-nodal continuous piecewise-affine surfaces
\(q_d^U,q_d^L\in\mathcal F_d\) with \(A_dq_d^U=A_dq_d^L=m_d\), and rational
numbers
\[
 \underline U_d\leq U_d(m_d)\leq\overline U_d,
 \qquad
 \underline L_d\leq L_d(m_d)\leq\overline L_d,
\]
such that \(\overline U_d-\underline U_d\leq\epsilon\) and
\(\overline L_d-\underline L_d\leq\epsilon\).  The certificate verifies the
exact continuous cell integrals, nodal clipping inequalities, squared Euclidean
gradient inequalities, an atomic clipping-aware partial-transport plan, and the
continuous remainder \(M_dh(1+\lVert\lambda_d\rVert_1)\).  Thus the returned
rational intervals are outward enclosures, not floating-residual claims; no
bit-complexity bound is asserted.
\end{theorem}
\begin{proof}
Fix an arm and suppress its subscript when this causes no ambiguity.  Write
\(J=J_d\), \(M=M_d\), \(A=A_d\), \(m=m_d\),
\(\mathcal F=\mathcal F_d\), and
\[
 T(g)=\int_B g\,d\nu.
\]
The proof uses dovetailed enumeration of exact rational witnesses and exact
algebraic dual certificates.  The geometric support assumptions make every
triangulation below finite; the cell-density and boundary-weight assumptions
make all of its integrals exactly algebraic.  The relative-interior condition
is used only to prove termination.  In particular, no ambient full-rank
condition is imposed: all local perturbations take place in
\(E_d=\operatorname{span}(K_d-K_d)\).  The argument requires neither overlap
of different cell laws nor positive boundary mass inside a cell; it uses only
their stated supports and the normalization \(\nu(B)=1\).

First consider exact primal certificates.  On a rational conforming
triangulation \(\mathcal T_h\), a continuous piecewise-affine function is
represented by its nodal vector \(z\).  Because the triangulation conforms to
the cells and to all density breaks, exact integration gives
\[
 Az=C_hz,                                                     \tag{1}
\]
where \(C_h\) is a rational matrix.  Indeed, on a triangle with rational
vertices, its area is rational and the integral of each barycentric coordinate
is one third of that area; multiplication by the rational constant density
preserves rationality.  The nodal clipping restrictions are the finite rational
inequalities \(0\leq z_v\leq1\).  If \(R_T\) is the rational edge matrix of a
nondegenerate triangle \(T\), then
\[
 \nabla(z|_T)=R_T^{-\top}\Delta_Tz,
 \qquad
 \|\nabla(z|_T)\|_2^2\leq M^2.                              \tag{2}
\]
For rational \(z\) and rational \(M\), (2) is an exactly decidable rational
inequality.  Convexity of each triangle shows that nodal clipping implies
\(0\leq z(x)\leq1\) throughout the triangle, while (2) is exactly its
Euclidean Lipschitz restriction.  Thus (1)--(2) certify membership in
\(\mathcal F_{d,h}\), and hence in \(\mathcal F_d\), together with the exact
continuous moment equality \(Az=m\).

The boundary objective is also exactly representable.  On a conforming
boundary subsegment \(e\) of length \(\ell_e\), let \(z_0,z_1\) and
\(w_0,w_1\) be the endpoint values of the affine function and weight.
Parametrizing the segment by \(t\in[0,1]\) and integrating the product of the
two affine interpolants yields
\[
 \int_e zw\,d\mathcal H^1
 =\ell_e\,
   \frac{(2w_0+w_1)z_0+(w_0+2w_1)z_1}{6}.                  \tag{3}
\]
Every squared segment length is rational.  Hence for rational nodal data the
finite sum in (3) lies in a finite multiquadratic extension of \(\mathbb Q\),
and integer-arithmetic interval bisection gives rational lower and upper
enclosures of \(T(z)\) of arbitrarily small width.

We next prove that exact rational primal certificates are dense at each
endpoint.  Since every \(\mu_{dj}\) is a probability measure, the constant
surface \(1/2\) has moment vector
\[
 c=A(1/2)=\tfrac12\mathbf 1\in K_d.                         \tag{4}
\]
Suppose first that \(M>0\).  If \(m\neq c\), choose a sufficiently small
rational \(s\in(0,1)\) and define
\[
 m'=\frac{m-sc}{1-s}
    =m+\frac{s}{1-s}(m-c).                                  \tag{5}
\]
Because \(m-c\in E_d\) and \(r_d(m)>0\), choosing
\(s\|m-c\|_2/(1-s)<r_d(m)\) puts \(m'\) in the relative ball from
\(\mathrm{def:relative-interior-radius}\); indeed a smaller relative ball
about \(m'\) is still contained in that ball, so \(m'\in\operatorname{ri}K_d\).
By \(\mathrm{thm:relative-interior-programs}\), on all sufficiently fine
refinements there is \(q'\in\mathcal F_{d,h}\) with \(Aq'=m'\).  Therefore
\[
 q^0=(1-s)q'+s/2,
 \qquad Aq^0=m,                                              \tag{6}
\]
has all nodal values in \([s/2,1-s/2]\) and all triangle-gradient norms at
most \((1-s)M<M\).  If \(m=c\), the constant \(q^0=1/2\) supplies the same
strict range and gradient slack.

Let \(\eta>0\).  The convergence of the exact-moment inner programs in
\(\mathrm{thm:relative-interior-programs}\) supplies, on a sufficiently fine
mesh, a function \(p\in\mathcal F_{d,h}\) such that
\[
 Ap=m,
 \qquad T(p)>U_d(m)-\eta/4.                                 \tag{7}
\]
Pass \(p\) and \(q^0\) to a common member of the nested midpoint refinement.
For a sufficiently small positive mixing coefficient \(\delta\),
\[
 p^0=(1-\delta)p+\delta q^0,
 \qquad Ap^0=m,
 \qquad T(p^0)>U_d(m)-\eta/2,                               \tag{8}
\]
and every clipping and gradient inequality for \(p^0\) is strict.  The affine
solution set \(\{z:C_hz=m\}\) is defined over \(\mathbb Q\), since both
\(C_h\) and the assumed input \(m\) are rational.  Rational Gaussian
elimination provides a rational particular solution and a rational basis of
its nullspace.  Approximating the corresponding free coordinates of \(p^0\)
by rationals therefore produces rational points in the same exact affine
space converging to \(p^0\).  The strict inequalities in (8) persist under a
sufficiently small perturbation, and (3) is continuous.  Consequently the
enumeration eventually encounters a rational-nodal surface \(q^U\) satisfying
\[
 q^U\in\mathcal F_d,
 \qquad Aq^U=m,
 \qquad T(q^U)>U_d(m)-\eta.                                 \tag{9}
\]
The first two assertions in (9) are exactly certified by (1)--(2); the last is
used only in the mathematical proof that the search terminates.

If \(M=0\), convexity of \(S\) makes every member of \(\mathcal F_d\)
constant.  When \(J\neq\varnothing\), compatibility forces all entries of
\(m\) to equal one rational constant, so both endpoints and their constant
witness coincide and are exact.  When \(J=\varnothing\), there are no moment
constraints and the constants zero and one give \(L_d=0\) and \(U_d=1\).
Thus the asserted terminating construction is immediate in either zero-radius
case.  Henceforth take \(M>0\).

We now construct a certified dual upper enclosure.  By
\(\mathrm{thm:clipping-transport-dual}\), the condition \(r_d(m)>0\) implies
that a finite optimal unrestricted multiplier
\(\lambda^*\in\mathbb R^{|J|}\) exists.  The dual objective is continuous.
More precisely, compatibility and the range restriction give
\(0\leq m_j\leq1\) and \(0\leq\int g\,d\mu_{dj}\leq1\), whence, for any
\(\lambda,\lambda'\),
\begin{align}
 |D_d(\lambda;m)-D_d(\lambda';m)|
 &\leq |(\lambda-\lambda')^\top m|
   +\sup_{g\in\mathcal F_d}
       \left|\sum_{j\in J}(\lambda_j-\lambda'_j)
                    \int g\,d\mu_{dj}\right| \\
 &\leq2\|\lambda-\lambda'\|_1.                            \tag{10}
\end{align}
Thus unrestricted rational multipliers contain points with dual objective
arbitrarily close to \(U_d(m)\); no multiplier-sum restriction appears.

For a rational multiplier \(\lambda\), put
\[
 \sigma(\lambda)=\nu-\sum_{j\in J}\lambda_j\mu_{dj},
 \qquad
 \sigma_h(\lambda)=\nu_h-\sum_{j\in J}\lambda_j\mu_{dj,h}.\tag{11}
\]
The nearest-node Voronoi boundaries are rational affine lines because a
difference of squared distances to two rational nodes is rational affine.
Intersecting them with rational density triangles and rational boundary
segments therefore yields finitely many rational polygons and subsegments.
It follows that the cell-atom masses are rational and the boundary-atom masses
belong to the same kind of finite multiquadratic extension as (3).  All atom
masses of \(\sigma_h(\lambda)\) are consequently exact algebraic numbers.

For completeness, their comparisons require no floating decision.  Reduce the
finitely many square roots to generators whose square classes are independent
over \(\mathbb F_2\); products of distinct generators form a basis, so equality
is equality of rational coefficient vectors.  For a nonzero element, nested
rational interval evaluation eventually excludes zero.  This gives a
terminating exact sign test and hence the exact atomic Jordan decomposition
\(\sigma_h^+-\sigma_h^-\).

Enumerate finite matrices \(\pi=(\pi_{xy})\) with nonnegative rational entries
and accept exactly those satisfying
\[
 \sum_y\pi_{xy}\leq\sigma_h^+(\{x\}),
 \qquad
 \sum_x\pi_{xy}\leq\sigma_h^-(\{y\}).                     \tag{12}
\]
Every comparison in (12) terminates by the preceding sign procedure.  For an
accepted matrix, the attained partial-transport identity in
\(\mathrm{thm:clipping-transport-dual}\) implies
\[
 H_M(\sigma_h)
 \leq \mathcal C_h(\lambda,\pi)
 :=\sigma_h^+(S)-\sum_{x,y}\pi_{xy}
      +M\sum_{x,y}\|x-y\|_2\pi_{xy}.                       \tag{13}
\]
The right side is exact algebraic and admits arbitrarily tight rational outward
enclosures.  These rational plans approach the atomic minimum: start from an
attained optimal real plan, and approximate every positive entry from below by
rationals.  Row and column sums can only decrease, so (12) remains true, while
the finite expression in (13) converges to the optimum.  This also handles
zero capacities exactly.

It remains to pass from the atomic support value to the continuous one.  The
projection-distance conclusions of \(\mathrm{thm:relative-interior-programs}\),
the probability normalizations, and the \(M\)-Lipschitz restriction give, for
every \(g\in\mathcal F_d\),
\begin{align}
 &\left|\int g\,d\sigma(\lambda)-
             \int g\,d\sigma_h(\lambda)\right| \\
 &\quad\leq
 \left|\int_B\{g(x)-g(p_{B,h}(x))\}\,d\nu(x)\right|
 +\sum_{j\in J}|\lambda_j|
   \left|\int\{g(x)-g(p_{dj,h}(x))\}\,d\mu_{dj}(x)\right| \\
 &\quad\leq Mh(1+\|\lambda\|_1).                           \tag{14}
\end{align}
Taking suprema in (14), combining the result with (13), and using the dual
identity proves that every rational upward rounding of
\[
 \lambda^\top m+\mathcal C_h(\lambda,\pi)
       +Mh(1+\|\lambda\|_1)                                \tag{15}
\]
is a valid upper bound on \(U_d(m)\).  Formula (14) is the advertised
continuous remainder; it is an analytic inequality rather than a floating
mesh residual.  A rational downward rounding of (3) for any certified
\(q^U\) is, by feasibility, a valid lower bound on \(U_d(m)\).

Define the implementation of \(\mathsf{VPA}\) to dovetail over the midpoint
refinement level, bounds on numerators and denominators of rational nodal
vectors, unrestricted rational multipliers and rational plan entries, and the
precision of every algebraic isolating interval.  At each finite stage it
checks (1), the nodal clipping inequalities, (2), (12), and the rational
ordering of the proposed endpoint bounds.  Each individual check terminates.
It stops the upper-endpoint search as soon as the resulting rational enclosure
has width at most the requested rational \(\epsilon\).

To prove that this stopping event occurs, fix \(\eta>0\).  Equation (9)
supplies an enumerated primal candidate whose objective exceeds
\(U_d(m)-\eta\).  By (10), rational-multiplier density supplies a rational
\(\lambda\) with
\[
 D_d(\lambda;m)<U_d(m)+\eta.                                \tag{16}
\]
Keep this finite multiplier fixed.  Along the nested refinements,
\(h\downarrow0\), so (14) is eventually smaller than \(\eta\).  At such a
fixed finite atomic problem, the rational-plan approximation following (13)
supplies \(\pi\) whose cost is less than \(H_M(\sigma_h)+\eta\).  A second use
of (14) gives \(H_M(\sigma_h)\leq H_M(\sigma)+\eta\).  Consequently (15),
before outward rounding, is less than
\[
 D_d(\lambda;m)+3\eta<U_d(m)+4\eta.                         \tag{17}
\]
Choose each final upward and downward rounding error smaller than \(\eta\).
The certified lower endpoint then exceeds \(U_d(m)-2\eta\), and the certified
upper endpoint is below \(U_d(m)+5\eta\); their gap is less than \(7\eta\).
Taking \(7\eta<\epsilon\) proves that the dovetailed search encounters a
stopping certificate.  Neither \(r_d(m)\) nor a convergence rate is an input
to the search; relative interior proves existence of a finite certificate.

Reflection supplies the lower endpoint.  The bijection \(g\mapsto1-g\))
preserves clipping and the Lipschitz radius and gives
\[
 A_d(1-g)=\mathbf1-A_dg,
 \qquad
 L_d(m)=1-U_d(\mathbf1-m).                                  \tag{18}
\]
It maps \(K_d\) onto \(\mathbf1-K_d\), maps its tangent space isometrically,
and therefore preserves the relative-interior radius.  Apply the preceding
algorithm at the rational mean \(\mathbf1-m\).  If
\([\underline U',\overline U']\) encloses
\(U_d(\mathbf1-m)\), then
\[
 [1-\overline U',\,1-\underline U']                         \tag{19}
\]
encloses \(L_d(m)\) with the same width.  If \(q'\) is the exact primal
witness produced in the reflected problem, then \(q^L=1-q'\) is rational
nodal, belongs to \(\mathcal F_d\), and satisfies \(A_dq^L=m\).  Dovetailing
the two endpoint searches and the two arms returns all four enclosures and
witnesses, and halts once each requested gap is at most \(\epsilon\).
Equations (1)--(3), (12)--(15), and (18) are exactly the continuous-moment,
clipping, squared-gradient, atomic partial-transport, and remainder checks
listed in the theorem.  No bit-complexity claim has been used.

Finally, the following exact two-cell instance verifies constructively that
the search is not relying on asymptotic floating residuals.  Take
\[
 S=[-1,1]\times[0,1],\qquad B=\{0\}\times[0,1],
\]
let \(C_1=[0,1]\times[0,1/2]\) and
\(C_2=[0,1]\times[1/2,1]\), give each cell its uniform probability law, let
\(\nu\) be uniform on \(B\), and set \(M=1/2\) and
\(m=(1/2,1/2)\).  For \(j=1,2\), put \(y_1=1/4\), \(y_2=3/4\), and define the
explicit rational piecewise-affine hats
\[
 \psi_j(x,y)=\max\{0,\,1-8|x-1/2|-8|y-y_j|\}.
\]
The support of \(\psi_j\) is a rational diamond strictly inside \(C_j\).
Consequently \(\int\psi_j\,d\mu_k=0\) for \(k\neq j\), whereas
\(a_j:=\int\psi_j\,d\mu_j>0\) is rational by triangle integration.  For all
sufficiently small \((u_1,u_2)\), the surface
\[
 g_u=\frac12+\frac{u_1}{a_1}\psi_1+
                  \frac{u_2}{a_2}\psi_2
\]
retains strict clipping and has Lipschitz constant below \(1/2\), while its
moment vector is \(m+(u_1,u_2)\).  Thus a Euclidean neighborhood of \(m\)
lies in \(K_d\), proving \(r_d(m)>0\) without a rank assumption beyond the
displayed two independent moments.

The mixture \(\rho=(\mu_1+\mu_2)/2\) is uniform on
\([0,1]\times[0,1]\).  Coupling \((x,y)\) to \((0,y)\) and using the
Lipschitz inequality gives every feasible \(g\) with these moments the bound
\[
 T(g)-\frac12\sum_{j=1}^2\int g\,d\mu_j
 \leq M\int_0^1\!\int_0^1x\,dx\,dy=\frac14.                \tag{20}
\]
The rational piecewise-affine surfaces
\[
 q^U(x,y)=\frac34-\frac12\max\{x,0\},
 \qquad
 q^L(x,y)=\frac14+\frac12\max\{x,0\}                       \tag{21}
\]
have both cell means exactly \(1/2\) and boundary means \(3/4\) and \(1/4\),
respectively.  Equation (20) and reflection therefore prove exactly
\[
 U_d(m)=\frac34,
 \qquad L_d(m)=\frac14.                                    \tag{22}
\]
For the rational multiplier \(\lambda=(1/2,1/2)\), let \(\bar h_k\) be a
rational outward upper bound for the mesh diameter with
\(\bar h_k\downarrow0\).  Push the coupling in (20) through the boundary and
cell projections.  Each transport distance rises by at most
\(2\bar h_k\).  In this axis-aligned dyadic instance all atomic capacities are
rational; after canceling common source and target mass, a rational finite
flow feasible for (12) therefore certifies
\[
 H_M\!\left(\nu_h-\tfrac12\mu_{1,h}-\tfrac12\mu_{2,h}\right)
 \leq\frac14+\bar h_k.                                     \tag{23}
\]
Here \(\lambda^\top m=1/2\), while the continuous remainder (14) is
\(\bar h_k\).  The exact witnesses (21) and the outward dual certificate (23)
give
\[
 \frac34\leq U_d(m)\leq\frac34+2\bar h_k,
 \qquad
 \frac14-2\bar h_k\leq L_d(m)\leq\frac14.                 \tag{24}
\]
These verified rational gaps shrink to zero and every primal moment in (21)
is an exact continuous cell integral.  This completes the construction and
the proof.
\end{proof}

\begin{lemma}[lem:piccoli-rossi-flat-duality]\label{lem:piccoli-rossi-flat-duality}
Let \(\alpha\) and \(\beta\) be finite nonnegative Radon measures on Euclidean
space, and define
\[
 d_{\mathrm{flat}}(\alpha,\beta)
 =\sup\left\{\int \phi\,d(\alpha-\beta):
       \lVert\phi\rVert_\infty\leq1,\ \operatorname{Lip}(\phi)\leq1\right\}.
\]
Then Piccoli and Rossi's generalized distance satisfies
\(W_1^{1,1}(\alpha,\beta)=d_{\mathrm{flat}}(\alpha,\beta)\).
\end{lemma}

\begin{lemma}[lem:piccoli-rossi-submeasure-attainment]\label{lem:piccoli-rossi-submeasure-attainment}
For finite nonnegative Radon measures \(\alpha\) and \(\beta\), the value
\(W_1^{1,1}(\alpha,\beta)\) equals
\[
 \min_{\substack{\widetilde\alpha\leq\alpha,\ \,
                  \widetilde\beta\leq\beta\\
                  \lvert\widetilde\alpha\rvert=\lvert\widetilde\beta\rvert}}
 \left\{\lvert\alpha-\widetilde\alpha\rvert
       +\lvert\beta-\widetilde\beta\rvert
       +W_1(\widetilde\alpha,\widetilde\beta)\right\},
\]
and the minimizing submeasures are attained.
\end{lemma}

\begin{lemma}[lem:hoeffding-bounded-sum]\label{lem:hoeffding-bounded-sum}
If \(X_1,\ldots,X_n\) are independent and satisfy
\(a_i\leq X_i\leq b_i\) almost surely, then, for every \(t>0\),
\[
 \Pr\left\{\frac1n\sum_{i=1}^n(X_i-\mathbb EX_i)\geq t\right\}
 \leq \exp\left\{-\frac{2n^2t^2}{\sum_{i=1}^n(b_i-a_i)^2}\right\}.
\]
\end{lemma}

\begin{theorem}[thm:affine-hull-projected-confidence-inversion]\label{thm:affine-hull-projected-confidence-inversion}
Let \(Q_n=\widetilde Q_{0,n}\times\widetilde Q_{1,n}\) be the affine-hull-projected confidence region in \(\mathrm{def:confidence-box}\), and let \(C_n\) be its whole-set inversion in \(\mathrm{def:whole-set-confidence}\).  Under conditional cell-sample independence,
\[
 \inf_{\mathbb P}\Pr_{\mathbb P}\{I(m_0,m_1)\subseteq C_n\}\geq1-\gamma,
\]
where the infimum ranges over all maintained \([0,1]\)-valued cell-outcome laws having the fixed score laws and compatible means.  Whenever \(Q_n\cap(K_0\times K_1)\neq\varnothing\), \(C_n\) is an attained compact interval equal to the effect image of all response-surface pairs whose joint moment vector lies in \(Q_n\).  If \(r_d(m_d)>0\) for both arms and \(\min_{d,j}n_{dj}\to\infty\), then
\[
 \Pr\{Q_n\cap(K_0\times K_1)\neq\varnothing\}\longrightarrow1,
 \qquad
 d_{\mathrm H}\bigl(C_n,I(m_0,m_1)\bigr)\xrightarrow{\mathbb P}0.
\]
The consistency conclusion requires no condition \(E_d=\mathbb R^{|J_d|}\).
\end{theorem}
\begin{proof}
\paragraph{Causal object, observed information, and projected region.}
Fix a maintained sampling law \(\mathbb P\) and condition on the positive, nonrandom cell counts \(n_{dj}\).  Let \(k_d=|J_d|\), let \(\widehat m_d=(\widehat m_{dj})_{j\in J_d}\), and put
\[
 \mathcal A_d=\operatorname{aff}K_d,
 \qquad
 E_d=\operatorname{span}(K_d-K_d).
\]
Compatibility in \(\mathrm{ass:mean-compatibility}\) gives \(m_d\in K_d\subseteq\mathcal A_d\), so \(\mathcal A_d=m_d+E_d\) is a nonempty closed affine subspace of \(\mathbb R^{k_d}\).  The fixed score laws in \(\mathrm{ass:known-cell-score-laws}\) determine \(A_d\) and \(K_d\); the only randomness considered here is that of the cell outcomes.

The causal target is \(\tau_B=\int_B(f_1-f_0)\,d\nu\).  It is distinct from the random confidence inversion.  By \(\mathrm{thm:sharp-mean-information-set}\), the sharp mean-information set generated by the population moments is
\[
 I(m_0,m_1)
 =\bigl[L_1(m_1)-U_0(m_0),\ U_1(m_1)-L_0(m_0)\bigr].       \tag{1}
\]
The observed-law region \(Q_n\) and its inversion \(C_n\) are defined in \(\mathrm{def:confidence-box}\) and \(\mathrm{def:whole-set-confidence}\); they do not redefine the causal target or the set in (1).

We first record the projection inequality used below.  If \(\mathcal A=a+E\) is a nonempty affine subspace of a finite-dimensional Euclidean space, its orthogonal projection satisfies
\[
 P_{\mathcal A}x-P_{\mathcal A}y=P_E(x-y).
\]
The orthogonal decomposition of \(x-y\) therefore gives
\[
 \|P_{\mathcal A}x-P_{\mathcal A}y\|_2^2
 \leq
 \|P_E(x-y)\|_2^2+\|(I-P_E)(x-y)\|_2^2
 =\|x-y\|_2^2.                                             \tag{2}
\]
Thus projection is nonexpansive.  Since \(m_d\in\mathcal A_d\), (2) and the definition of \(\widetilde m_d\) imply
\[
 \|\widetilde m_d-m_d\|_2
 =\|P_{\mathcal A_d}\widehat m_d-P_{\mathcal A_d}m_d\|_2
 \leq\|\widehat m_d-m_d\|_2.                              \tag{3}
\]
This is an elementary orthogonal-decomposition argument, not a rank or inversion assumption.

No cross-boundary overlap probability is invoked or divided by in this argument.  The boundary and assignment restrictions enter through the already established sharp-set theorem, while the only boundary functional fact used directly below is the normalization \(\nu(B)=1\).

\paragraph{Simultaneous finite-sample coverage.}
Suppose first that \(N\geq1\).  For every observed cell \((d,j)\), \(\mathrm{ass:cell-sample-independence}\), the declared support \(0\leq Y_{dji}\leq1\), and \(\mathrm{lem:hoeffding-bounded-sum}\) give
\[
 \Pr_{\mathbb P}\{\widehat m_{dj}-m_{dj}>e_{dj}\}
 \leq \exp\{-2n_{dj}e_{dj}^{2}\}.
\]
Applying the same cited lemma to \(-Y_{dji}\in[-1,0]\) gives the other tail.  Since \(n_{dj}\geq1\), \(0<\gamma<1\), and
\[
 e_{dj}^{2}=\frac{\log(2N/\gamma)}{2n_{dj}},
\]
we obtain
\[
 \Pr_{\mathbb P}\{|\widehat m_{dj}-m_{dj}|>e_{dj}\}
 \leq2e^{-2n_{dj}e_{dj}^{2}}
 =\frac{\gamma}{N}.                                       \tag{4}
\]
Finite subadditivity over the \(N\) cells yields the event
\[
 \mathcal E_n
 =\bigcap_{d\in\{0,1\}}\bigcap_{j\in J_d}
   \{|\widehat m_{dj}-m_{dj}|\leq e_{dj}\},
 \qquad
 \Pr_{\mathbb P}(\mathcal E_n)\geq1-\gamma.               \tag{5}
\]
On \(\mathcal E_n\), equations (3) and (5) give, arm by arm,
\[
 \|\widetilde m_d-m_d\|_2
 \leq\|\widehat m_d-m_d\|_2
 \leq\left(\sum_{j\in J_d}e_{dj}^{2}\right)^{1/2}
 =\rho_{d,n}.                                               \tag{6}
\]
Because \(m_d\in\mathcal A_d\), (6) is exactly the membership \(m_d\in\widetilde Q_{d,n}\).  Compatibility then implies
\[
 (m_0,m_1)\in Q_n\cap(K_0\times K_1).                      \tag{7}
\]
The definition of \(C_n\) and (7) show that \(I(m_0,m_1)\subseteq C_n\) on \(\mathcal E_n\).  Consequently
\[
 \Pr_{\mathbb P}\{I(m_0,m_1)\subseteq C_n\}
 \geq\Pr_{\mathbb P}(\mathcal E_n)
 \geq1-\gamma.                                             \tag{8}
\]
No step from (2) through (8) depends on the particular maintained law beyond the declared independence, support, cell means, and fixed score laws.  Taking the infimum over those laws proves the finite-sample assertion.  If \(N=0\), all intersections, sums, and products over cells are empty; the moment spaces are zero-dimensional, orthogonal projection fixes their unique point, and (7)--(8) hold with probability one.  Thus this edge case requires no evaluation of a nonexistent cell radius.

\paragraph{Exact random effect image and attained endpoints.}
By \(\mathrm{ass:polygonal-score-support}\), \(S\) is compact.  The class \(\mathcal F_d\) in \(\mathrm{def:surface-class}\) is uniformly bounded and has the common modulus
\[
 |g(x)-g(y)|\leq M_d\|x-y\|_2.                              \tag{9}
\]
To verify compactness without an undeclared theorem, take finite \(1/\ell\)-nets of \(S\), successively extract subsequences convergent on each net, and diagonalize.  If \(x\) is within \(1/\ell\) of a net point \(x_\ell\), (9) gives
\[
 |g_r(x)-g_s(x)|
 \leq 2M_d/\ell+|g_r(x_\ell)-g_s(x_\ell)|.                \tag{10}
\]
First increasing \(\ell\) and then increasing \(r,s\) shows that the diagonal subsequence is uniformly Cauchy.  Its uniform limit remains \([0,1]\)-valued and \(M_d\)-Lipschitz, so it lies in \(\mathcal F_d\).  Hence \(\mathcal F_d\) is compact in the uniform norm.  Its defining range and Lipschitz inequalities also show that it is convex.

Every \(\mu_{dj}\) is a probability measure, and \(\nu(B)=1\) by \(\mathrm{ass:boundary-weight-normalization}\).  Therefore
\[
 \left|\int_S(g-\bar g)\,d\mu_{dj}\right|
 \leq\|g-\bar g\|_\infty,
 \qquad
 \left|\int_B(g-\bar g)\,d\nu\right|
 \leq\|g-\bar g\|_\infty.                                \tag{11}
\]
Thus the moment and boundary maps are continuous.  In particular, \(K_d=A_d(\mathcal F_d)\) is compact and convex.

Each \(\widetilde Q_{d,n}\) is the intersection of a closed Euclidean ball with the finite-dimensional affine subspace \(\mathcal A_d\), hence is compact and convex.  Define
\[
 \mathcal R_n
 =\{(g_0,g_1)\in\mathcal F_0\times\mathcal F_1:
        (A_0g_0,A_1g_1)\in Q_n\}
\]
and
\[
 T(g_0,g_1)=\int_B(g_1-g_0)\,d\nu.
\]
By (11), \(\mathcal R_n\) is a closed convex subset of the compact convex product \(\mathcal F_0\times\mathcal F_1\).  Moreover,
\[
 \mathcal R_n\neq\varnothing
 \quad\Longleftrightarrow\quad
 Q_n\cap(K_0\times K_1)\neq\varnothing,                   \tag{12}
\]
because \(K_d=A_d(\mathcal F_d)\) by \(\mathrm{def:moment-body}\).  When (12) holds, the continuous linear image \(T(\mathcal R_n)\) is a nonempty compact convex subset of \(\mathbb R\), and hence an interval whose endpoints are attained.  Partitioning \(\mathcal R_n\) by its joint moment vector and applying the fiberwise attainment and sharpness in \(\mathrm{thm:sharp-mean-information-set}\) gives
\[
\begin{aligned}
 T(\mathcal R_n)
 &=\bigcup_{(q_0,q_1)\in Q_n\cap(K_0\times K_1)}
   \left\{\int_B(g_1-g_0)\,d\nu:
       g_d\in\mathcal F_d,\ A_dg_d=q_d\right\}\\
 &=\bigcup_{(q_0,q_1)\in Q_n\cap(K_0\times K_1)}I(q_0,q_1)
 =C_n.                                                       \tag{13}
\end{aligned}
\]
This proves the exact effect-image and attainment claims.  In particular, \(C_n\) is an observed-law inversion of sharp causal fibers, not a pointwise regression interval.

\paragraph{Relative-affine endpoint modulus.}
Fix an arm and suppress its subscript temporarily.  Convexity of \(\mathcal F_d\), linearity of \(A_d\), and linearity of the boundary functional imply that \(U_d\) is concave on \(K_d\): convexly mixing maximizers at two moment vectors is feasible at the mixed moment and attains the mixed objective.  The same argument makes \(L_d\) convex.  The range restriction and \(\nu(B)=1\) yield
\[
 0\leq L_d(q)\leq U_d(q)\leq1
 \qquad(q\in K_d).                                         \tag{14}
\]
Consequently both \(U_d\) and \(1-L_d\) are concave functions taking values in \([0,1]\).

Assume \(r_d(m_d)>0\), and choose any fixed
\[
 0<R_d<r_d(m_d).                                            \tag{15}
\]
The feasible radii in \(\mathrm{def:relative-interior-radius}\) are downward closed.  Since \(R_d\) lies strictly below their supremum, there is a feasible radius larger than \(R_d\), and hence
\[
 m_d+R_d\mathbb B_{E_d}\subseteq K_d.                      \tag{16}
\]
Take \(q\in K_d\), put \(t=\|q-m_d\|_2\), and suppose first that \(t>0\).  Since \(q-m_d\in E_d\), (16) implies
\[
 z_-=m_d-\frac{R_d}{t}(q-m_d)\in K_d,
 \qquad
 m_d=\frac{R_d}{R_d+t}q+\frac{t}{R_d+t}z_-.                \tag{17}
\]
Concavity and nonnegativity of \(U_d\), followed by its upper bound in (14), give
\[
 U_d(m_d)
 \geq\frac{R_d}{R_d+t}U_d(q)
       +\frac{t}{R_d+t}U_d(z_-)
 \geq\frac{R_d}{R_d+t}U_d(q),
\]
and therefore
\[
 U_d(q)-U_d(m_d)
 \leq\frac{t}{R_d}U_d(m_d)
 \leq\frac{t}{R_d}.                                       \tag{18}
\]
If \(t<R_d\), then, for every \(u\in E_d\) with \(\|u\|_2\leq R_d-t\),
\[
 \|q+u-m_d\|_2\leq t+(R_d-t)=R_d,
\]
so (16) implies \(q+(R_d-t)\mathbb B_{E_d}\subseteq K_d\).  Applying (18) with center \(q\), comparison point \(m_d\), and radius \(R_d-t\) gives
\[
 U_d(m_d)-U_d(q)\leq\frac{t}{R_d-t}.                       \tag{19}
\]
Equations (18)--(19), including the trivial case \(t=0\), show that
\[
 |U_d(q)-U_d(m_d)|\leq\frac{2t}{R_d}
 \quad\text{whenever }q\in K_d\text{ and }t\leq R_d/2.    \tag{20}
\]
Repeating the same argument for the concave \([0,1]\)-valued function \(1-L_d\) yields
\[
 |L_d(q)-L_d(m_d)|\leq\frac{2t}{R_d}
 \quad\text{under the same condition}.                    \tag{21}
\]
All displacement vectors in (16)--(21) lie in \(E_d\).  Thus these bounds use relative-interior slack only and nowhere require \(E_d=\mathbb R^{k_d}\).

\paragraph{Eventual nonemptiness and Hausdorff consistency.}
Keep \(\gamma\in(0,1)\), the finite cell sets, and the maintained score laws fixed, and suppose \(\min_{d,j}n_{dj}\to\infty\).  Because a random variable in \([0,1]\) has variance at most \(1/4\), conditional independence gives
\[
 \mathbb E_{\mathbb P}\|\widehat m_d-m_d\|_2^2
 =\sum_{j\in J_d}\operatorname{Var}_{\mathbb P}(\widehat m_{dj})
 \leq\sum_{j\in J_d}\frac{1}{4n_{dj}}
 \longrightarrow0.                                        \tag{22}
\]
For every \(c>0\), the pointwise inequality
\(\mathbf 1\{\|\widehat m_d-m_d\|_2>c\}
 \leq\|\widehat m_d-m_d\|_2^2/c^2\)
and (22) prove \(\widehat m_d-m_d\xrightarrow{\mathbb P}0\).  Equation (3) then gives
\[
 \|\widetilde m_d-m_d\|_2\xrightarrow{\mathbb P}0.        \tag{23}
\]
Moreover, directly from the definition of the projected radius,
\[
 \rho_{d,n}^{2}
 =\sum_{j\in J_d}\frac{\log(2N/\gamma)}{2n_{dj}}
 \longrightarrow0.                                        \tag{24}
\]
Define
\[
 \Delta_{d,n}=\|\widetilde m_d-m_d\|_2+\rho_{d,n}.
\]
Equations (23)--(24) imply
\[
 \Delta_{d,n}\xrightarrow{\mathbb P}0.                   \tag{25}
\]

On the event \(\|\widetilde m_d-m_d\|_2<R_d\), the facts that \(\widetilde m_d\in\mathcal A_d=m_d+E_d\) and (16) imply \(\widetilde m_d\in K_d\).  The center of the projected ball belongs to that ball, so
\[
 \widetilde Q_{d,n}\cap K_d\neq\varnothing.               \tag{26}
\]
By (23), (26) holds for both arms with probability tending to one, proving eventual nonemptiness of the joint inversion.

For any \(q_d\in\widetilde Q_{d,n}\cap K_d\), the triangle inequality and the definition of \(\widetilde Q_{d,n}\) give
\[
 \|q_d-m_d\|_2
 \leq\|q_d-\widetilde m_d\|_2+\|\widetilde m_d-m_d\|_2
 \leq\rho_{d,n}+\|\widetilde m_d-m_d\|_2
 =\Delta_{d,n}.                                             \tag{27}
\]
Write the population interval in (1) as \([a,b]\).  On the event \(\Delta_{d,n}\leq R_d/2\) for both arms, (20), (21), and (27) imply, uniformly over \((q_0,q_1)\in Q_n\cap(K_0\times K_1)\),
\[
\begin{aligned}
 |\{L_1(q_1)-U_0(q_0)\}-a|
 &\leq \frac{2\Delta_{1,n}}{R_1}+\frac{2\Delta_{0,n}}{R_0},\\
 |\{U_1(q_1)-L_0(q_0)\}-b|
 &\leq \frac{2\Delta_{1,n}}{R_1}+\frac{2\Delta_{0,n}}{R_0}. \tag{28}
\end{aligned}
\]
By (13), the lower endpoint of \(C_n\) is the minimum of the first expression in braces in (28), and the upper endpoint is the maximum of the second.  For nonempty compact intervals, Hausdorff distance is the maximum absolute difference of corresponding endpoints.  Therefore (25), (26), and (28) imply
\[
 d_{\mathrm H}\bigl(C_n,I(m_0,m_1)\bigr)
 \leq\frac{2\Delta_{1,n}}{R_1}+\frac{2\Delta_{0,n}}{R_0}
 \xrightarrow{\mathbb P}0.                                \tag{29}
\]
Assigning any value to the Hausdorff distance on the event that \(C_n\) is empty does not change (29), because that event has probability tending to zero by (26).  This completes the finite-sample coverage, exact-image, attainment, eventual-nonemptiness, and relative-affine consistency arguments.
\end{proof}

\section*{Technical internal limitation (diagnostic only)}

At a boundary mean, finite multiplier attainment and exact piecewise-affine inner feasibility can fail even though the continuous primal endpoint exists. A unit-rectangle zero mean forces the true surface to vanish on the observed cell while finite dual slopes leave a positive cone value; a separate multi-cell distance-function example has a legal exposed-face mean with no exactly feasible function in any finite piecewise-affine triangulation. These are regression tests, not claimed contributions against a published estimator.

\section{Honest scope}

The population identification, clipping-aware transport representation, planar formulas, relative-interior approximation bounds, affine-hull-projected whole-set inference, and terminating verified reconstruction routine are proved contributions. Reconstruction is guaranteed for rational polygonal geometry, rational piecewise-constant cell densities, rational piecewise-affine boundary weights, rational inputs, and relative-interior means; no bit-complexity rate is claimed. At boundary-face means, finite multiplier attainment and exact finite-mesh primal feasibility can fail, so no universal boundary-face reconstruction guarantee is made. Known within-cell score laws and separately selected arm-specific radii remain maintained nonstandard inputs, not facts established by Qiu and Stoye or for Dell's data. The response class is exactly the \([0,1]\)-bounded isotropic Euclidean-Lipschitz class applied armwise, while the fixed-areal operator, target, inference, and certification are the new objects. The application therefore reports the full prespecified score-law and radius sensitivity frontier and certifies a sign only when it is invariant across that grid. Sharpness conditions only on cell means, not full within-cell outcome distributions; no floating-residual certificate, shrinking-cell asymptotic, or empirical reanalysis is claimed.

\begin{thebibliography}{99}

  \bibitem{QiuStoye2026BoundaryDecision} Qiu, Chen, and Jörg Stoye. 2026. Statistical Decisions and Partial Identification: With Application to Boundary Discontinuity Design. Advances in Economics and Econometrics, Thirteenth World Congress, forthcoming. arXiv:2601.17648.

  \bibitem{CattaneoTitiunikYu2025AggregationBoundary} Cattaneo, Matias D., Rocío Titiunik, and Ruiqi Rae Yu. 2025. Causal Treatment Effect Aggregation in Boundary Discontinuity Designs. Working paper, July 4, 2025.

  \bibitem{CattaneoTitiunikYu2025LocationBoundary} Cattaneo, Matias D., Rocío Titiunik, and Ruiqi Rae Yu. 2025. Estimation and Inference in Boundary Discontinuity Designs. arXiv:2505.05670.

  \bibitem{PiccoliRossi2016GeneralizedWasserstein} Piccoli, Benedetto, and Francesco Rossi. 2016. On Properties of the Generalized Wasserstein Distance. Archive for Rational Mechanics and Analysis 222:1339--1365. doi:10.1007/s00205-016-1026-7.

  \bibitem{MulaNouy2024MomentSOSOT} Mula, Olga, and Anthony Nouy. 2024. Moment--SoS Methods for Optimal Transport Problems. Numerische Mathematik 156:1541--1578. doi:10.1007/s00211-024-01422-x.

  \bibitem{BaldiMourrain2025ExactMoments} Baldi, Lorenzo, and Bernard Mourrain. 2025. Exact Moment Representation in Polynomial Optimization. Journal of Symbolic Computation 128:102403. doi:10.1016/j.jsc.2024.102403.

  \bibitem{BaldiMourrain2023EffectivePutinar} Baldi, Lorenzo, and Bernard Mourrain. 2023. On the Effective Putinar Positivstellensatz and Moment Approximation. Mathematical Programming 200:71--103. doi:10.1007/s10107-022-01877-6.

  \bibitem{Moon2024AggregatePartialID} Moon, Sarah. 2024. Partial Identification of Individual-Level Parameters Using Aggregate Data in a Nonparametric Model. arXiv:2403.07236, revised 2025.

  \bibitem{BeresteanuMolchanovMolinari2011RandomSets} Beresteanu, Arie, Ilya Molchanov, and Francesca Molinari. 2011. Sharp Identification Regions in Models with Convex Moment Predictions. Econometrica 79:1785--1821. doi:10.3982/ECTA8680.

  \bibitem{KaidoMolinariStoye2022ConstraintQualifications} Kaido, Hiroaki, Francesca Molinari, and Jörg Stoye. 2022. Constraint Qualifications in Partial Identification. Econometric Theory 38:1215--1233. doi:10.1017/S0266466621000207.

  \bibitem{ChernozhukovHongTamer2007SetInference} Chernozhukov, Victor, Han Hong, and Elie Tamer. 2007. Estimation and Confidence Regions for Parameter Sets in Econometric Models. Econometrica 75:1243--1284. doi:10.1111/j.1468-0262.2007.00794.x.

  \bibitem{ArmstrongKolesar2020HonestCI} Armstrong, Timothy B., and Michal Kolesár. 2020. Simple and Honest Confidence Intervals in Nonparametric Regression. Quantitative Economics 11:1--39. doi:10.3982/QE1199.

  \bibitem{DeanerKwon2025ComonotoneRD} Deaner, Ben, and Soonwoo Kwon. 2025. Extrapolation in Regression Discontinuity Design Using Comonotonicity. arXiv:2507.00289.

  \bibitem{Hoeffding1963ProbabilityInequalities} Hoeffding, Wassily. 1963. Probability Inequalities for Sums of Bounded Random Variables. Journal of the American Statistical Association 58:13--30.

  \bibitem{Dell2010Mita} Dell, Melissa. 2010. The Persistent Effects of Peru's Mining Mita. Econometrica 78:1863--1903. doi:10.3982/ECTA8121.

\end{thebibliography}

\end{document}